\chapter{同伦 \texorpdfstring{$n$}{n}类型}
\label{cha:hlevels}

\index{n-type@$n$-type|(}%
\indexsee{h-level}{$n$-type}

同伦论的一个基本概念是\emph{同伦 $n$-类型（homotopy $n$-type）}：一个在维度 $n$ 以上不含任何有趣同伦的空间。
例如，同伦 $0$-类型本质上就是集合，不包含非平凡的道路；而同伦 $1$-类型可能包含非平凡的道路，但不包含道路之间的非平凡道路。
同伦 $n$-类型也称为\emph{$n$-截断空间（$n$-truncated spaces）}。
我们已经在 \cref{sec:basics-sets} 中提到过这个概念；本章的第一个目标就是在同伦类型论中给出它的精确定义。

与截断性对偶的概念是连通性：如果一个空间在维度 $n$ 及\emph{以下}不含任何有趣的同伦，则称其为\emph{$n$-连通的（$n$-connected）}。
例如，一个空间是 $0$-连通的（也简称为"连通的"），当且仅当它只有一个连通分支；是 $1$-连通的（也称为"单连通的"），当且仅当它也没有非平凡环路（尽管它可能有环路之间的非平凡高阶环路\index{loop!n-@$n$-}）。

将这两个概念都扩展到映射上，便能最清楚地看到截断性与连通性之间的对偶关系。
如果映射的所有纤维都是 $n$-截断的（或 $n$-连通的），我们就称该映射是\emph{$n$-截断的}或\emph{$n$-连通的}。
那么，$n$-连通映射和 $n$-截断映射就构成了一个\emph{正交分解系统}
\index{orthogonal factorization system}
\indexsee{factorization!system, orthogonal}{orthogonal factorization system}
中的两个映射类，也就是说，每个映射都可以唯一地分解为一个 $n$-连通映射后接一个 $n$-截断映射。

在 $n={-1}$ 的情形，$n$-截断映射就是嵌入，而 $n$-连通映射就是满射，正如在 \cref{sec:mono-surj} 中所定义的那样。
因此，$n$-连通分解系统是集合间函数标准像分解（即分解为满射后接单射）的大规模推广。
在本章末尾，我们会简要概述一个更为一般的理论：任何类型论的\emph{模态}都会引出一个类似的分解系统。

\section{\texorpdfstring{$n$}{n}类型的定义}
\label{sec:n-types}

正如在 \cref{sec:basics-sets,sec:contractibility} 中提到的，从零以下两个层级开始定义 $n$-类型是比较方便的：$(-1)$-类型就是命题，$(-2)$-类型则是可缩的类型。

\begin{defn}\label{def:hlevel}
  对于 $n \geq -2$，我们通过递归定义谓词 $\istype{n} : \type \to \type$ 如下：
  \[ \istype{n}(X) \defeq
  \begin{cases}
    \iscontr(X) & \text{ if } n = -2, \\
    \prd{x,y : X} \istype{n'}(\id[X]{x}{y}) & \text{ if } n = n'+1.
  \end{cases}
  \]
  如果 $\istype{n}(X)$ 有实例，我们就说 $X$ 是一个\define{$n$-类型（$n$-type）}，有时也说它是\emph{$n$-截断的}，
  \indexdef{n-type@$n$-type}%
  \indexsee{n-truncated@$n$-truncated!type}{$n$-type}%
  \indexsee{type!n-type@$n$-type}{$n$-type}%
  \indexsee{type!n-truncated@$n$-truncated}{$n$-type}。%
\end{defn}

\begin{rmk}
  \cref{def:hlevel} 中的数字 $n$ 取遍所有大于等于 $-2$ 的整数。
  我们可以形式化地理解这一点：可以定义一个这样的整数的类型 $\Z_{{\geq}-2}$（其归纳原理与 $\nat$ 相同），或者对每个 $k : \nat$ 定义一个谓词 $\istype{(k-2)}$。
  无论哪种方式，我们都可以通过对 $n$ 进行归纳来证明关于 $n$-类型的定理，其中 $n = -2$ 是基础情形。
\end{rmk}

\begin{eg}
  \index{set}
  我们在 \cref{thm:prop-minusonetype} 中看到，$X$ 是 $(-1)$-类型当且仅当它是一个命题。
  因此，$X$ 是 $0$-类型当且仅当它是一个集合。
\end{eg}

我们也已经看到存在不是集合的类型（\cref{thm:type-is-not-a-set}）。
然而，到目前为止，我们还没有对任何 $n>0$ 证明存在不是 $n$-类型的类型。
不过在 \cref{cha:homotopy} 中，我们将证明 $(n+1)$-球面 $\Sn^{n+1}$ 不是 $n$-类型。
（Kraus（克劳斯）也证明了第 $n$ 层嵌套的泛等宇宙也不是 $n$-类型，且未使用任何高阶归纳类型。）
此外，在 \cref{sec:whitehead} 中，我们将给出一个对\emph{任何}（有限）数字 $n$ 都不是 $n$-类型的例子。

我们先证明 $n$-类型在某些运算和构造子下封闭，以此开始 $n$-类型的一般理论。

\begin{thm}\label{thm:h-level-retracts}
  \index{retract!of a type}%
  \index{retraction}%
 设 $p : X \to Y$ 是一个收缩。对任意 $n\geq -2$，若 $X$ 是 $n$-类型，则 $Y$ 也是 $n$-类型。
\end{thm}

\begin{proof}
 我们对 $n$ 进行归纳。
 基础情形 $n=-2$ 由 \cref{thm:retract-contr} 处理。

 对于归纳步骤，假设任何 $n$-类型的收缩核仍是 $n$-类型，并且 $X$ 是一个 $\nplusone$-类型。
 任取 $y, y' : Y$；我们需要证明 $\id{y}{y'}$ 是一个 $n$-类型。
 令 $s$ 为 $p$ 的一个截面，并令 $\epsilon$ 为一个同伦 $\epsilon : p \circ s \htpy 1$。
 由于 $X$ 是一个 $\nplusone$-类型，$\id[X]{s(y)}{s(y')}$ 是一个 $n$-类型。
 我们断言 $\id{y}{y'}$ 是 $\id[X]{s(y)}{s(y')}$ 的一个收缩核。
 对于截面，我们取
 \[ \apfunc s : (y=y') \to (s(y)=s(y')). \]
 对于收缩，我们定义 $t:(s(y)=s(y'))\to(y=y')$ 为
 \[ t(q) \defeq  \opp{\epsilon_y} \ct \ap p q \ct \epsilon_{y'}.\]
 为了证明 $t$ 是 $\apfunc s$ 的一个收缩，我们需要证明对于任意 $r:y=y'$，有
 \[ \opp{\epsilon_y} \ct \ap p {\ap sr} \ct \epsilon_{y'} = r \]
 但这由 \cref{lem:htpy-natural} 可得。
\end{proof}

作为直接推论，我们得到 $n$-类型在等价下的稳定性（这从泛等性也能直接推出）：

\begin{cor}\label{cor:preservation-hlevels-weq}
 如果 $\eqv{X}{Y}$ 且 $X$ 是一个 $n$-类型，那么 $Y$ 也是。
\end{cor}

此外，回忆 \cref{sec:mono-surj} 中引入的嵌入概念。

\begin{thm}\label{thm:isntype-mono}
  \index{function!embedding}
  若 $f:X\to Y$ 是一个嵌入，且对于某个 $n\ge -1$，$Y$ 是一个 $n$-类型，那么 $X$ 也是。
\end{thm}

\begin{proof}
  任取 $x,x':X$；我们需要证明 $\id[X]{x}{x'}$ 是一个 $\nminusone$-类型。
  但由于 $f$ 是嵌入，我们有 $(\id[X]{x}{x'}) \eqvsym (\id[Y]{f(x)}{f(x')})$，而由假设后者是一个 $\nminusone$-类型。
\end{proof}

注意，这一定理在 $n=-2$ 时不成立：映射 $\emptyt \to \unit$ 是一个嵌入，但 $\unit$ 是一个 $(-2)$-类型而 $\emptyt$ 不是。

\begin{thm}\label{thm:hlevel-cumulative}
 $n$-类型的层级在如下意义下是累积的：
   给定一个数 $n \geq -2$，如果 $X$ 是一个 $n$-类型，那么它也是一个 $\nplusone$-类型。
\end{thm}

\begin{proof}
 我们对 $n$ 进行归纳。

 当 $n = -2$ 时，我们需要证明一个可缩类型，比如 $A$，具有可缩的道路空间。
       设 $a_0: A$ 是 $A$ 的收缩中心，并设 $x, y : A$。我们证明 $\id[A]{x}{y}$
       是可缩的。
       由 $A$ 的可缩性，我们有一条道路 $\contr_x \ct \opp{\contr_y} : x = y$，我们选择它作为
       $\id{x}{y}$ 的收缩中心。
       给定任意 $p : x = y$，我们需要证明 $p = \contr_x \ct \opp{\contr_y}$。
           根据道路归纳，只需证明
       $\refl{x} = \contr_x \ct \opp{\contr_x}$，而这是平凡的。

 对于归纳步骤，我们需要证明：若 $X$ 是 $\nplusone$-类型，则
           $\id[X]{x}{y}$ 是 $\nplusone$-类型。将归纳假设应用于 $\id[X]{x}{y}$
           即得所需结论。
\end{proof}

% \section{Preservation under constructors}
% \label{sec:ntype-pres}

我们现在证明，大多数类型形成运算都保持 $n$-类型。

\begin{thm}\label{thm:ntypes-sigma}
 设 $n \geq -2$，且令 $A : \type$ 与 $B : A \to \type$。
 如果 $A$ 是一个 $n$-类型，并且对于所有 $a : A$，$B(a)$ 都是 $n$-类型，那么
 $\sm{x : A} B(x)$ 也是 $n$-类型。
\end{thm}

\begin{proof}
 我们对 $n$ 进行归纳。

 当 $n = -2$ 时，我们选择 $\sm{x : A} B(x)$ 的收缩中心为偶对
       $(a_0, b_0)$，其中 $a_0 : A$ 是 $A$ 的收缩中心，而 $b_0 : B(a_0)$ 是 $B(a_0)$ 的收缩中心。
       给定 $\sm{x : A} B(x)$ 中任意其他元素 $(a,b)$，我们分别利用 $A$ 和 $B(a_0)$ 的可缩性
       提供一条道路 $\id{(a, b)}{(a_0,b_0)}$。

 对于归纳步骤，假设 $A$ 是 $\nplusone$-类型，并且
         对于任意 $a : A$，$B(a)$ 是 $\nplusone$-类型。我们证明 $\sm{x : A} B(x)$ 是 $\nplusone$-类型：
      固定 $\sm{x : A} B(x)$ 中的 $(a_1, b_1)$ 和 $(a_2,b_2)$，
     我们证明 $\id{(a_1, b_1)}{(a_2,b_2)}$ 是一个 $n$-类型。
      根据 \cref{thm:path-sigma}，我们有
      \[ \eqvspaced{(\id{(a_1, b_1)}{(a_2,b_2)})}{\sm{p : \id{a_1}{a_2}} (\id[B(a_2)]{\trans{p}{b_1}}{b_2})} \]
   并且由 $n$-类型在等价下的保持性（\cref{cor:preservation-hlevels-weq}），
   只需证明后者是 $n$-类型。这由归纳假设可得。
\end{proof}

作为一个特例，如果 $A$ 和 $B$ 是 $n$-类型，那么 $A\times B$ 也是。
另请注意，\cref{thm:hlevel-cumulative} 蕴含：如果 $A$ 是 $n$-类型，那么对于任意 $x,y:A$，$\id[A]xy$ 也是 $n$-类型。
将此与 \cref{thm:ntypes-sigma} 结合，我们看到对于任意 $n$-类型之间的函数 $f:A\to C$ 和 $g:B\to C$，它们的拉回\index{pullback}
\[ A\times_C B \defeq \sm{x:A}{y:B} (f(x)=g(y)) \]
（参见 \cref{ex:pullback}）也是 $n$-类型。
更一般地，$n$-类型在所有 \emph{极限} 下封闭。

\begin{thm}\label{thm:hlevel-prod}
 设 $n\geq -2$，且令 $A : \type$ 与 $B : A \to \type$。
 如果对于所有 $a : A$，$B(a)$ 是 $n$-类型，那么 $\prd{x : A} B(x)$ 也是 $n$-类型。
\end{thm}

\begin{proof}
  我们对 $n$ 进行归纳。
  当 $n = -2$ 时，结果直接就是 \cref{thm:contr-forall}。

  对于归纳步骤，假设结论对 $n$-类型成立，并且每个 $B(a)$ 是 $\nplusone$-类型。
  令 $f, g : \prd{a:A}B(a)$。
  我们需要证明 $\id{f}{g}$ 是 $n$-类型。
  根据函数外延性以及 $n$-类型在等价下的封闭性，只需证明 $\prd{a : A} (\id[B(a)]{f(a)}{g(a)})$ 是 $n$-类型。
  这由归纳假设可得。
\end{proof}

作为上述定理的一个特例，如果 $B$ 是一个 $n$-类型，那么函数空间 $A \to B$ 也是 $n$-类型。
现在我们可以推广 \cref{cha:basics} 中的观察，即 $\isset(A)$ 和 $\isprop(A)$ 是命题：

\begin{thm}\label{thm:isaprop-isofhlevel}
 对于任意 $n \geq -2$ 和任意类型 $X$，类型 $\istype{n}(X)$ 都是一个命题。
\end{thm}

\begin{proof}
  我们对 $n$ 进行归纳。

 对于基情形，需要证明对于任意 $X$，类型 $\iscontr(X)$ 是一个命题。
 这正是 \cref{thm:isprop-iscontr}。

对于归纳步骤，我们需要证明
\[\prd{X : \type} \isprop (\istype{n}(X)) \to \prd{X : \type} \isprop (\istype{\nplusone}(X)). \]
为了证明此蕴含关系的结论，我们需要证明：对于任意类型 $X$，类型
\[\prd{x, x' : X}\istype{n}(x = x')\]
是一个命题。根据 \cref{thm:isprop-forall} 或 \cref{thm:hlevel-prod}，只需证明对于任意 $x, x' : X$，类型 $\istype{n}(x =_X x')$ 是一个命题。
而这由对类型 $(x =_X x')$ 施加归纳假设即可得出。
\end{proof}

最后，我们证明 $n$-类型构成的类型本身是一个 $\nplusone$-类型。
我们将其定义为：
\symlabel{universe-of-ntypes}
\[\ntype{n} \defeq \sm{X : \type} \istype{n}(X). \]
必要时，我们可以写作 $\ntypeU{n}$ 以指定宇宙 $\UU$。
特别地，如 \cref{cha:basics} 中所定义，我们有 $\prop \defeq \ntype{(-1)}$ 与 $\set \defeq \ntype{0}$。
需要注意的是，正如对于 \prop 和 \set 一样，由于 $\istype{n}(X)$ 是一个命题，根据 \cref{thm:path-subset}，对于任何 $(X,p), (X',p'):\ntype{n}$ 我们有
\begin{align*}
  \Big(\id[\ntype{n}]{(X, p)}{(X', p')}\Big) &\eqvsym (\id[\type] X X')\\
  &\eqvsym (\eqv{X}{X'}).
\end{align*}

\begin{thm}\label{thm:hleveln-of-hlevelSn}
 对于任意 $n \geq -2$，类型 $\ntype{n}$ 是一个 $\nplusone$-类型。
\end{thm}

\begin{proof}%[Proof of \cref{thm:hleveln-of-hlevelSn}]
  令 $(X, p), (X', p') : \ntype{n}$；我们需要证明 $\id{(X, p)}{(X', p')}$ 是一个 $n$-类型。
  根据上述观察，这个类型等价于 $\eqv{X}{X'}$。
  接下来，我们注意到投影
  \[(\eqv{X}{X'}) \to (X \rightarrow X').\]
  是一个嵌入，因此如果 $n\geq -1$，那么由 \cref{thm:isntype-mono}，只需证明 $X \rightarrow X'$ 是一个 $n$-类型。
  但由于 $n$-类型在箭头类型下是封闭的，这就化归到 $X'$ 是 $n$-类型这一假设。

  在 $n=-2$ 的情形下，上述论证表明 $\eqv{X}{X'}$ 是一个 $(-1)$-类型——但它也是有实例的，因为任何两个可缩的类型都等价于 \unit，从而彼此等价。
  因此，$\eqv{X}{X'}$ 也是一个 $(-2)$-类型。
\end{proof}

\section{恒等证明唯一性与 Hedberg 定理}
\label{sec:hedberg}

\index{set|(}%

在 \cref{sec:basics-sets} 中，我们定义了什么是\emph{集合（set）}：如果对于任意 $x, y : X$ 以及 $p, q : x =_X y$ 都有 $p = q$，则称类型 $X$ 是一个集合。
在传统类型理论中，这一性质被称为\define{恒等证明唯一性（uniqueness of identity proofs, UIP）}。
\indexdef{uniqueness!of identity proofs}%
我们也已经看到，这等价于它是上一节意义下的 $0$-类型。
下面给出另一种等价的刻画，涉及 Streicher 的``K 公理'' \cite{Streicher93}：

\begin{thm}\label{thm:h-set-uip-K}
 一个类型 $X$ 是集合，当且仅当它满足\define{K 公理（Axiom K）}：
 \indexdef{axiom!Streicher's Axiom K}%
 即，对于所有 $x : X$ 和 $p : (x =_A x)$，都有 $p = \refl{x}$。
\end{thm}

\begin{proof}
 显然，K 公理是 UIP 的一个特例。
 反之，若 $X$ 满足 K 公理，任取 $x, y : X$ 和 $p, q : (\id{x}{y})$；我们要证明 $p=q$。
 但对 $q$ 进行归纳恰好将这一目标归约为 K 公理。
\end{proof}

我们要强调，\emph{我们}并没有将 UIP 或 K 原则作为公理来假设！
它们只是某些类型可能具有、也可能不具有的性质（即等价于该类型为集合）。
回顾 \cref{thm:type-is-not-a-set}，\emph{并非}所有类型都是集合。

下面的定理提供了另一种证明类型是集合的有用方法。

\begin{thm}\label{thm:h-set-refrel-in-paths-sets}
  \index{relation!reflexive}%
  设 $R$ 是类型 $X$ 上的一个自反的\index{reflexivity!of a relation}纯粹关系，且蕴含恒等。
  那么 $X$ 是一个集合，并且对于所有 $x,y:X$，$R(x,y)$ 都等价于 $\id[X]{x}{y}$。
\end{thm}

\begin{proof}
  令 $\rho : \prd{x:X} R(x,x)$ 见证 $R$ 的自反性，并令 \narrowequation{f : \prd{x,y:X} R(x,y) \to (\id[X]{x}{y})} 见证 $R$ 蕴含恒等。
  我们首先注意到，定理中的两个陈述是等价的。
  因为一方面，如果 $X$ 是一个集合，那么 $\id[X]xy$ 就是一个命题，而根据假设，它与命题 $R(x,y)$ 逻辑等价，因此也必然与 $R(x,y)$ 等价。
  另一方面，如果 $\id[X]xy$ 等价于 $R(x,y)$，那么与后者一样，它对所有 $x,y:X$ 都是命题，因此 $X$ 是一个集合。

  我们给出这个定理的两种证明。
  第一种直接证明 $X$ 是集合；第二种直接证明 $R(x,y)\eqvsym (x=y)$。

  \emph{第一种证明：}我们证明 $X$ 是集合。
  思路与 \cref{thm:prop-set} 相同：函数 $f$ 对其参数 $x$ 和 $y$ 必须是"连续的"。
  不过，由于我们需要处理类型为 $R(x,y)$ 的额外参数，符号上会稍微复杂一些。

  首先，对任意 $x:X$ 和 $p:\id[X]xx$，考虑 $\apdfunc{f(x)}(p)$。
  这是从 $f(x,x)$ 到自身的一条依赖道路。
  由于 $f(x,x)$ 仍然是一个函数 $R(x,x) \to (\id[X]xx)$，根据 \cref{thm:dpath-arrow}，对任意 $r:R(x,x)$ 我们得到一条道路
  \[\trans{p}{f(x,x,r)} = f(x,x,\trans{p}r).
  \]
  在等式的左边，是恒等类型中的传输，也就是拼接。
  而在右边，由于 $r$ 和 $\trans{p}r$ 都位于命题 $R(x,x)$ 中，我们有 $\trans{p}r = r$。
  因此，将 $r$ 替换为 $\rho(x)$，我们得到
  \[ f(x,x,\rho(x)) \ct p = f(x,x,\rho(x)). \]
  通过消去，得到 $p=\refl{x}$。
  所以 $X$ 满足 K 公理，从而是集合。

  \emph{第二种证明：}我们证明每个 $f(x,y) : R(x,y) \to \id[X]{x}{y}$ 都是一个等价。
  根据 \cref{thm:total-fiber-equiv}，只需证明 $f$ 诱导了一个全空间之间的等价：
  \begin{equation*}
    \eqv{\Parens{\sm{y:X}R(x,y)}}{\Parens{\sm{y:X}\id[X]{x}{y}}}.
  \end{equation*}
  由 \cref{thm:contr-paths}，右边的类型是可缩的，因此我们只需证明左边的类型也是可缩的。
  我们取偶对 $\pairr{x,\rho(x)}$ 作为收缩的中心。
  余下只需证明，对每个 ${y:X}$ 和每个 ${H:R(x,y)}$，都有
  \begin{equation*}
    \id{\pairr{x,\rho(x)}}{\pairr{y,H}}.
  \end{equation*}
  但由于 $R(x,y)$ 是一个命题，根据 \cref{thm:path-sigma}，我们只需证明 $\id[X]{x}{y}$，而这可以由 $f(H)$ 给出。
\end{proof}

\begin{cor}\label{notnotstable-equality-to-set}
  如果一个类型 $X$ 满足以下性质：对任意 $x,y:X$，$\neg\neg(x=y)\to(x=y)$，那么 $X$ 是一个集合。
\end{cor}

另一种证明类型是集合的便捷方法如下。
回顾 \cref{sec:intuitionism}，一个类型 $X$ 被称为具有\emph{可判定相等性（decidable equality）}
\index{decidable!equality|(}%
如果对所有 $x, y : X$，我们有
\[(x =_X y) + \neg (x =_X y).\]
\index{continuity of functions in type theory@``continuity'' of functions in type theory}%
\index{functoriality of functions in type theory@``functoriality'' of functions in type theory}%
这是一个非常强的条件：它意味着，当道路 $x=y$ 存在时，可以关于 $x$ 和 $y$ "连续地"（或"可计算地"或"函子性地"）选取该道路。
结合 \cref{thm:h-set-refrel-in-paths-sets} 与下面的引理，这蕴含 $X$ 是一个集合。

\begin{lem}\label{lem:hedberg-helper}
对于任意类型 $A$，我们有 $(A+\neg A)\to(\neg\neg A\to A)$。
\end{lem}

\begin{proof}
这在 \cref{thm:not-lem} 中已经基本证明过了，但我们再重复一下论证过程。
假设 $x:A+\neg A$。我们需要考虑两种情况。
如果 $x$ 是某个 $a:A$ 的左注入 $\inl(a)$，那么我们有一个常值函数 $\neg\neg A
\to A$，它将所有元素都映射到 $a$。如果 $x$ 是某个 $t:\neg A$ 的右注入 $\inr(t)$，
那么对于每个 $g:\neg\neg A$，都有 $g(t):\emptyt$。因此我们可以使用
\emph{ex falso quodlibet}，即 $\rec{\emptyt}$，来为任意 $g:\neg\neg A$ 获得一个 $A$ 中的元素。
\end{proof}

\index{anger}


\begin{thm}[Hedberg]\label{thm:hedberg}
  \index{Hedberg's theorem}%
  \index{theorem!Hedberg's}%
  如果 $X$ 有可判定相等性，那么 $X$ 是一个集合。
\end{thm}

\begin{proof}
如果 $X$ 有可判定相等性，那么对于任意 $x,y:X$，$\neg\neg(x=y)\to(x=y)$ 成立。因此，Hedberg 定理可由
\cref{notnotstable-equality-to-set} 推出。
\end{proof}

当然，这个定理与 \cref{thm:not-lem} 有密切的联系。
被 \cref{thm:not-lem} 所否定的命题 \LEM{\infty}，显然蕴含了每个类型都具有可判定相等性，从而都是集合——然而我们知道实际情况并非如此。
\index{excluded middle}%
请注意，从 \cref{sec:intuitionism} 引入的一致公理 \LEM{}，仅蕴含每个类型具有\emph{命题意义下的可判定相等性}，即对于任意 $A$，我们有
\indexdef{equality!merely decidable}%
\indexdef{merely!decidable equality}%
\[ \prd{a,b:A} (\brck{a=b} + \neg\brck{a=b}). \]

\index{decidable!equality|)}%

作为 \cref{thm:hedberg} 的一个应用示例，回顾 \cref{thm:nat-set}，我们曾利用 \cref{sec:compute-nat} 中对自然数相等类型的刻画，观察到 $\nat$ 是一个集合。
一个更传统的证明该定理的方法，仅使用 \eqref{eq:zero-not-succ} 和 \eqref{eq:suc-injective}，而不是 \cref{thm:path-nat} 的完整刻画，并用 \cref{thm:hedberg} 来填补空白。

\begin{thm}\label{prop:nat-is-set}
 自然数类型 $\nat$ 具有可判定相等性，因此是一个集合。
\end{thm}

\begin{proof}
  给定 $x, y : \nat$，我们通过对 $x$ 作归纳、对 $y$ 作分情形讨论来证明 $(x=y)+\neg(x=y)$。
  若 $x \jdeq 0$ 且 $y \jdeq 0$，我们取 $\inl(\refl0)$。
  若 $x \jdeq 0$ 且 $y \jdeq \suc(n)$，则由 \eqref{eq:zero-not-succ} 可得 $\neg (0 = \suc (n))$。

  对于归纳步骤，设 $x \jdeq \suc (n)$。
  若 $y \jdeq 0$，我们再次使用 \eqref{eq:zero-not-succ}。
  最后，若 $y \jdeq \suc (m)$，归纳假设给出 $(m = n)+\neg(m = n)$。
  在第一种情形中，若 $p:m=n$，则 $\ap \suc p:\suc(m)=\suc(n)$。
  而在第二种情形中，\eqref{eq:suc-injective} 推出 $\neg(\suc(m)=\suc(n))$。
\end{proof}

\index{set|)}%

\index{axiom!Streicher's Axiom K!generalization to n-types@generalization to $n$-types}%
尽管 Hedberg 定理似乎专属于集合（$0$-类型），但 K 公理可以自然地推广到 $n$-类型。
注意，普通的 K 公理（作为类型 $X$ 的一个性质）断言：对所有 $x:X$，环路空间 $\Omega(X,x)$（参见 \cref{def:loopspace}）是可缩的。
由于 $\Omega(X,x)$ 总是有实例的（由 $\refl{x}$ 给出），这等价于它是一个命题（一个 $(-1)$-类型）。
因为 $0 = (-1)+1$，这启发我们如下推广。

\begin{thm}\label{thm:hlevel-loops}
  对于任意 $n\geq -1$，类型 $X$ 是一个 $\nplusone$-类型，当且仅当对所有 $x : X$，类型 $\Omega(X, x)$ 是一个 $n$-类型。
\end{thm}

在证明此定理前，我们先证明一个辅助引理：

\begin{lem}\label{lem:hlevel-if-inhab-hlevel}
  给定 $n \geq -1$ 和 $X : \type$。
  如果已知 $X$ 有实例便能推出 $X$ 是 $n$-类型，那么 $X$ 本身就是 $n$-类型。
\end{lem}

\begin{proof}
  令 $f : X \to \istype{n}(X)$ 为给定的映射。
  我们需要证明，对任意 $x, x' : X$，类型 $\id{x}{x'}$ 是一个 $\nminusone$-类型。
  而此时 $f(x)$ 表明 $X$ 是一个 $n$-类型，因而其所有道路空间都是 $\nminusone$-类型。
\end{proof}

\begin{proof}[Proof of \cref{thm:hlevel-loops}]
  “仅当”方向是显然的，因为 $\Omega(X,x)\defeq (\id[X]xx)$。
  反过来，为证 $X$ 是一个 $\nplusone$-类型，我们需要证明：对任意 $x, x' : X$，类型 $\id{x}{x'}$ 是一个 $n$-类型。
  根据 \cref{lem:hlevel-if-inhab-hlevel}，只需构造一个映射
  \[ (\id{x}{x'}) \to \istype{n}(\id{x}{x'}). \]
  由道路归纳，只需考虑 $x\jdeq x'$ 的情形，此时由 $\Omega(X, x)$ 是 $n$-类型的假设即得证。
\end{proof}

\index{whiskering}
通过归纳和一些略显巧妙的须接（whiskering）操作，我们可以将 K 性质推广到 $n>0$ 的情形。

\begin{thm}\label{thm:ntype-nloop}
  \index{loop space!iterated}%
  对于每个 $n\ge -1$，类型 $A$ 是 $n$-类型当且仅当对所有的 $a:A$，$\Omega^{n+1}(A,a)$ 都是可缩的。
\end{thm}

\begin{proof}
  回顾 $\Omega^0(A,a) = (A,a)$，$n=-1$ 的情形即为 \cref{ex:prop-inhabcontr}。
  $n=0$ 的情形是 \cref{thm:h-set-uip-K}。
  现在我们使用归纳法；假设命题对 $n:\N$ 成立。
  根据 \cref{thm:hlevel-loops}，$A$ 是 $(n+1)$-类型当且仅当对所有 $a:A$，$\Omega(A,a)$ 都是 $n$-类型。
  由归纳假设，后者等价于：对所有 $p:\Omega(A,a)$，$\Omega^{n+1}(\Omega(A,a),p)$ 都是可缩的。

  由于 $\Omega^{n+2}(A,a) \defeq \Omega^{n+1}(\Omega(A,a),\refl{a})$ 且 $\Omega^{n+1} = \Omega^n \circ \Omega$，只需证明在带点类型 $\pointed\type$ 中 $\Omega(\Omega(A,a),p)$ 与 $\Omega(\Omega(A,a),\refl{a})$ 相等。
  为此，只需给出一个等价
  \[ g : \Omega(\Omega(A,a),p) \eqvsym \Omega(\Omega(A,a),\refl{a}) \]
  使其将基点 $\refl{p}$ 映射到基点 $\refl{\refl{a}}$。
  对于 $q:p=p$，定义 $g(q):\refl{a} = \refl{a}$ 为如下复合：
  \[ \refl{a} = p\ct \opp p \overset{q}{=} p\ct\opp p = \refl{a}, \]
  其中标记为“$q$”的道路实际上是 $\apfunc{\lam{r} r\ct\opp p} (q)$。
  那么 $g$ 是一个等价，因为它是以下一系列等价的复合：
  \[ (p=p) \xrightarrow{\apfunc{\lam{r} r\ct\opp p}} (p\ct \opp p = p\ct \opp p) \xrightarrow{i\ct - \ct \opp i} (\refl{a} = \refl{a}). \]
  这里用到了 \cref{eg:concatequiv,thm:paths-respects-equiv}，其中 $i:\refl{a} = p\ct \opp p$ 是典范的相等。
  显然有 $g(\refl{p}) = \refl{\refl{a}}$。
\end{proof}

\section{截断}
\label{sec:truncations}

\indexsee{n-truncation@$n$-truncation}{truncation}%
\index{truncation!n-truncation@$n$-truncation|(defstyle}%

在 \cref{subsec:prop-trunc} 中，我们引入了命题截断，它给出了一个类型到“仅仅是一个命题”（即 $(-1)$-类型）的“最佳逼近”。
在 \cref{sec:hittruncations} 中，我们将此截断构造为一个高阶归纳类型，并给出了一种将其推广到 0-截断的方法。
现在我们来解释一种更好的推广，它可以将任何类型截断为任意 $n\geq -2$ 的 $n$-类型；在经典同伦论中，这被称为它的 \define{$n$ 阶 Postnikov 截面（$n^{\mathrm{th}}$ Postnikov section）}。\index{Postnikov tower}

其思路是利用 \cref{thm:ntype-nloop} 与 \cref{lem:susp-loop-adj}。
前者指出，$A$ 是 $n$-类型当且仅当对所有 $a:A$，$\Omega^{n+1}(A,a)$
\index{loop space!iterated}%
是可缩的；后者则意味着
\narrowequation{\Omega^{n+1}(A,a) \eqvsym \Map_{*}(\Sn^{n+1},(A,a)),}
其中 $\Sn^{n+1}$ 配备了一个基点，我们不妨称之为 \base。
然而，$\Map_*(\Sn^{n+1},(A,a))$ 的可缩性是我们可以通过直接给出道路构造子来保证的。

\index{hub and spoke}%
我们将使用类似于 \cref{sec:hubs-spokes} 中的“毂辐构造（hub and spoke）”。
因此，对 $n\ge -1$，我们取 $\trunc nA$ 为由以下生成子生成的高阶归纳类型：

\begin{itemize}
\item 一个函数 $\tprojf n : A \to \trunc n A$，
\item 对于每个 $r:\Sn^{n+1} \to \trunc n A$，一个\emph{轮毂}点 $h(r):\trunc n A$，以及
\item 对于每个 $r:\Sn^{n+1} \to \trunc n A$ 和每个 $x:\Sn^{n+1}$，一条\emph{辐条}道路 $s_r(x):r(x) = h(r)$。
\end{itemize}

\noindent
有了这些构造子，就足以证明：

\begin{lem}
  $\trunc n A$ 是一个 $n$-类型。
\end{lem}

\begin{proof}
  根据 \cref{thm:ntype-nloop}，只需证明对所有 $b:\trunc nA$，$\Omega ^{n+1}(\trunc nA,b)$ 是可缩的。
  而由 \cref{lem:susp-loop-adj}，这等价于 \narrowequation{\Map_*(\Sn^{n+1},(\trunc nA,b)).}
  我们以后者为收缩中心：取值为 $b$ 的常值函数 $c_b:\Sn^{n+1} \to \trunc nA$，连同道路 $\refl b : c_b(\base) = b$。

  现在，$\Map_*(\Sn^{n+1},(\trunc nA,b))$ 中的任意元素由一个映射 $r:\Sn^{n+1} \to \trunc n A$ 和一条道路 $p:r(\base)=b$ 组成。
  根据函数外延性，要证明 $r = c_b$，只需对每个 $x:\Sn^{n+1}$ 给出一条道路 $r(x)=c_b(x) \jdeq b$。
  我们将其取为复合 $s_r(x) \ct \opp{s_r(\base)} \ct p$，其中 $s_r(x)$ 是在 $x$ 处的辐条。


  最后，我们必须证明，当沿着这个相等 $r=c_b$ 输运时，道路 $p$ 会变成 $\refl b$。
  根据道路类型中的输运，这意味着我们需要
  \[\opp{(s_r(\base) \ct \opp{s_r(\base)} \ct p)} \ct p = \refl b.\]
  而这直接由道路运算可得。
\end{proof}

（此构造在 $n=-2$ 时失败，但在那种情况下，我们可以对所有 $A$ 简单地定义 $\trunc{-2}{A}\defeq \unit$。
从现在起，我们假设 $n\ge -1$。）

\index{induction principle!for truncation}%
为了证明 $n$-截断所期望的泛性质，我们需要归纳原理。
我们像往常一样从构造子中提取它；该原理是说，给定 $P:\trunc nA\to\type$ 以及


\begin{itemize}
\item 对于每个 $a:A$，一个元素 $g(a) : P(\tproj na)$，
\item 对于每个 $r:\Sn^{n+1} \to \trunc n A$ 和 $r':\prd{x:\Sn^{n+1}} P(r(x))$，一个元素 $h'(r,r'):P(h(r))$，
\item 对于每个 $r:\Sn^{n+1} \to \trunc n A$ 和 $r':\prd{x:\Sn^{n+1}} P(r(x))$，以及每个 $x:\Sn^{n+1}$，一条依值道路
$\dpath{P}{s_r(x)}{r'(x)}{h'(r,r')}$，
\end{itemize}


则存在一个截面 $f:\prd{x:\trunc n A} P(x)$ 使得对所有 $a:A$ 有 $f(\tproj n a) \jdeq g(a)$。
为使其更为有用，我们将其重新表述如下。

\begin{thm}\label{thm:truncn-ind}
  对于任意类型族 $P:\trunc n A \to \type$，若每个 $P(x)$ 都是 $n$-类型，则对任意函数 $g : \prd{a:A} P(\tproj n a)$，存在一个截面 $f:\prd{x:\trunc n A} P(x)$ 使得对所有 $a:A$ 有 $f(\tproj n a)\defeq g(a)$。
\end{thm}

\begin{proof}
  只需构造上面列出的第二和第三项数据即可，因为 $g$ 的类型恰好就是第一项数据。
  给定 $r:\Sn^{n+1} \to \trunc n A$ 和 $r':\prd{x:\Sn^{n+1}} P(r(x))$，我们有 $h(r):\trunc n A$ 和 $s_r :\prd{x:\Sn^{n+1}} (r(x) =
h(r))$。
  定义 $t:\Sn^{n+1} \to P(h(r))$ 为 $t(x) \defeq \trans{s_r(x)}{r'(x)}$。
  由于 $P(h(r))$ 是 $n$-截断的，故存在一个点 $u:P(h(r))$ 和一个收缩 $v:\prd{x:\Sn^{n+1}} (t(x) = u)$。
  定义 $h'(r,r') \defeq u$，这给出了第二项数据。
  那么（回忆依值道路的定义），$v$ 的类型恰好就是第三项数据所要求的类型。
\end{proof}

特别地，如果 $E$ 是某个 $n$-类型，我们可以考虑在 $A$ 的每一点处取值恒为 $E$ 的常值类型族。
\symlabel{extend}
\index{recursion principle!for truncation}%
于是，每个映射 $f:A\to{}E$ 都可以延拓为一个映射 $\extend{f}:\trunc nA\to{}E$，其定义为 $\extend{f}(\tproj na)\defeq f(a)$；这就是 $\trunc n A$ 的\emph{递归原理}。

归纳原理还蕴含了这种形式的函数的唯一性原理。
\index{uniqueness!principle, propositional!for functions on a truncation}%
即，如果 $E$ 是 $n$-类型，且 $g,g':\trunc nA\to{}E$ 满足对每个 $a:A$ 有 $g(\tproj na)=g'(\tproj na)$，那么对所有 $x:\trunc nA$ 有 $g(x)=g'(x)$，因为类型 $g(x)=g'(x)$ 是 $n$-类型。
因此，$g=g'$。
（事实上，当 $E$ 是 $\nplusone$-类型时，此唯一性原理更一般地也成立。）
由此我们得到以下泛性质。

\begin{lem}[截断的泛性质]\label{thm:trunc-reflective}
  \index{universal!property!of truncation}%
  设 $n\ge-2$，$A:\type$ 且 $B:\typele{n}$。则以下映射是一个等价：
  \[\function{(\trunc nA\to{}B)}{(A\to{}B)}{g}{g\circ\tprojf n}\]
\end{lem}

\begin{proof}
  给定 $B$ 是 $n$-截断的，任何 $f:A\to{}B$ 都可以延拓为一个映射 $\extend{f}:\trunc nA\to{}B$。
  映射 $\extend{f}\circ\tprojf n$ 等于 $f$，因为根据定义，对于每个 $a:A$ 我们有 $\extend{f}(\tproj na)=f(a)$。
  而映射 $\extend{g\circ\tprojf n}$ 等于 $g$，因为它们都将 $\tproj na$ 映到 $g(\tproj na)$。
\end{proof}

用范畴论的语言来说，这表明 $n$-类型构成了类型范畴的一个\emph{反射子范畴}。
\index{reflective!subcategory}%
（要完全精确地陈述这一点，需要使用 $(\infty,1)$-范畴的语言。）
\index{.infinity1-category@$(\infty,1)$-category}%
特别地，这意味着 $n$-截断具有函子性：给定 $f:A\to B$，对复合 $A\xrightarrow{f} B \to \trunc n B$ 应用递归原理，就得到一个映射 $\trunc n f: \trunc n A \to \trunc n B$。
根据定义，我们有一个同伦
\begin{equation}
  \mathsf{nat}^f_n : \prd{a:A} \trunc n f(\tproj n a) = \tproj n {f(a)},\label{eq:trunc-nat}
\end{equation}
表达了映射 $\tprojf n$ 的\emph{自然性}。

唯一性蕴含了函子律，例如 $\trunc n {g\circ f} = \trunc n g \circ \trunc n f$ 和 $\trunc n{\idfunc[A]} = \idfunc[\trunc n A]$，以及相应的相干律。
我们还有高阶函子性，例如：

\begin{lem}\label{thm:trunc-htpy}
  给定 $f,g:A\to B$ 以及同伦 $h:f\htpy g$，存在一个诱导的同伦 $\trunc n h : \trunc n f \htpy \trunc n g$，使得复合
  \begin{equation}
    \xymatrix@C=3.6pc{\tproj n{f(a)} \ar@{=}[r]^-{\opp{\mathsf{nat}^f_n(a)}} &
      \trunc n f(\tproj n a) \ar@{=}[r]^-{\trunc n h(\tproj na)} &
      \trunc n g(\tproj n a) \ar@{=}[r]^-{\mathsf{nat}^g_n(a)} &
      \tproj n{g(a)}}\label{eq:trunc-htpy}
  \end{equation}
  等于 $\apfunc{\tprojf n}(h(a))$。
\end{lem}

\begin{proof}
  首先，我们确实有一个分量为 $\apfunc{\tprojf n}(h(a)) : \tproj n{f(a)} = \tproj n{g(a)}$ 的同伦。
  将其与路径 $\tproj n{f(a)} = \trunc n f(\tproj n a)$ 和 $\tproj n{g(a)} = \trunc n g(\tproj n a)$（分别源自 $\trunc n f$ 和 $\trunc ng$ 的定义）在两侧复合，便得到同伦 $(\trunc n f \circ \tprojf n) \htpy (\trunc n g \circ \tprojf n)$，进而由函数外延性得到一个相等。
  但由于 $(\blank\circ \tprojf n)$ 是一个等价，必然存在一条道路 $\trunc nf = \trunc ng$ 诱导出该同伦，而函数外延性的相干律蕴含了~\eqref{eq:trunc-htpy}。
\end{proof}

关于反射子范畴，以下观察也是标准的。

\begin{cor}
  一个类型 $A$ 是 $n$-类型，当且仅当 $\tprojf n : A \to \trunc n A$ 是一个等价。
\end{cor}

\begin{proof}
  “当”的部分由 $n$-类型在等价下封闭可得。
  反之，若 $A$ 是 $n$-类型，我们可以定义 $\ext(\idfunc[A]):\trunc n A\to{}A$。
  那么根据定义我们有 $\ext(\idfunc[A])\circ\tprojf n=\idfunc[A]:A\to{}A$。
  为了证明
  $\tprojf n\circ\ext(\idfunc[A])=\idfunc[\trunc nA]$，
  我们只需证明
  $\tprojf n\circ\ext(\idfunc[A])\circ\tprojf n=
  \idfunc[\trunc nA]\circ\tprojf n$。
  这一点同样成立：
  \[\raisebox{\depth-\height+1em}{\xymatrix{
    A \ar^-{\tprojf n}[r] \ar_{\idfunc[A]}[rd] &
    \trunc nA \ar^>>>{\ext(\idfunc[A])}[d] \ar@/^40pt/^{\idfunc[\trunc nA]}[dd] \\
    & A \ar_{\tprojf n}[d] \\
    & \trunc nA}}
  \qedhere\]
\end{proof}

$n$-类型的范畴还有一些并非所有反射子范畴都具有的特殊性质。
例如，反射子 $\trunc n-$ 保持有限积。

\begin{thm}\label{cor:trunc-prod}
  对于任意类型 $A$ 和 $B$，诱导的映射 $\trunc n{A\times B} \to \trunc nA \times \trunc nB$ 是一个等价。
\end{thm}

\begin{proof}
  只需证明 $\trunc nA \times \trunc nB$ 具有与 $\trunc n{A\times B}$ 相同的泛性质。
  于是，设 $C$ 是一个 $n$-类型；我们有如下等价：
  \begin{align*}
    (\trunc nA \times \trunc nB \to C)
    &= (\trunc nA \to (\trunc nB \to C))\\
    &= (\trunc nA \to (B \to C))\\
    &= (A \to (B \to C))\\
    &= (A \times B \to C)
  \end{align*}
  这里用到了 $\trunc nB$ 和 $\trunc nA$ 的泛性质，以及 $C$ 是 $n$-类型时 $B\to C$ 也是 $n$-类型这一事实。
  不难验证，此等价由与 $\tprojf n \times \tprojf n$ 复合给出，恰如所需。
\end{proof}

下面这个关于依赖和的事实也常常很有用。

\begin{thm}\label{thm:trunc-in-truncated-sigma}
设 $P:A\to\type$ 为一个类型族。则存在等价
\begin{equation*}
\eqv{\Trunc n{\sm{x:A}\trunc n{P(x)}}}{\Trunc n{\sm{x:A}P(x)}}.
\end{equation*}
\end{thm}

\begin{proof}
我们多次利用 $n$-截断的归纳原理，来构造函数
\begin{align*}
\varphi & : \Trunc n{\sm{x:A}\trunc n{P(x)}}\to\Trunc n{\sm{x:A}P(x)}\\
\psi & : \Trunc n{\sm{x:A}P(x)}\to \Trunc n{\sm{x:A} \trunc n{P(x)}}
\end{align*}
以及同伦 $H:\varphi\circ\psi\htpy \idfunc$ 与 $K:\psi\circ\varphi\htpy
\idfunc$，以表明它们互为拟逆。
定义 $\varphi$ 为 $\varphi(\tproj n{\pairr{x,\tproj nu}})\defeq\tproj n{\pairr{x,u}}$。
定义 $\psi$ 为 $\psi(\tproj n{\pairr{x,u}})\defeq\tproj n{\pairr{x,\tproj nu}}$。
然后定义 $H(\tproj n{\pairr{x,u}})\defeq \refl{\tproj n{\pairr{x,u}}}$ 与
$K(\tproj n{\pairr{x,\tproj nu}})\defeq \refl{\tproj n{\pairr{x,\tproj nu}}}$。
\end{proof}

\begin{cor}\label{thm:refl-over-ntype-base}
  若 $A$ 是一个 $n$-类型，且 $P:A\to\type$ 是任意类型族，则
  \[ \eqv{\sm{a:A} \trunc n{P(a)}}{\Trunc n{\sm{a:A}P(a)}} \]
\end{cor}

\begin{proof}
  如果 $A$ 是 $n$-类型，那么上式左边的类型本身就已经是 $n$-类型，因此等价于它的 $n$-截断；于是本推论由 \cref{thm:trunc-in-truncated-sigma} 可得。
\end{proof}

我们可以用与 \cref{sec:compute-coprod,sec:compute-nat} 中处理余积和自然数时相同的方法，来刻画截断的路径空间（该方法之后在 \cref{cha:homotopy} 计算同伦群时也会用到）。
不出所料，$A$ 的 $(n+1)$-截断中的路径空间，就是 $A$ 的路径空间的 $n$-截断。
确切地说，对于任意 $x,y:A$，存在一个典范映射
\begin{equation}
  f:\ttrunc n{x=_Ay}\to \Big(\tproj {n+1}x=_{\trunc{n+1}A}\tproj {n+1}y\Big)\label{eq:path-trunc-map}
\end{equation}
其定义为
\[f(\tproj n{p})\defeq \apfunc{\tprojf {n+1}}(p). \]
这个定义使用了 $\truncf n$ 的递归原理，这是可行的，因为 $\trunc {n+1}A$ 是 $(n+1)$-截断的，因此 $f$ 的余域是 $n$-截断的。

\begin{thm} \label{thm:path-truncation}
  对于任意 $A$、$x,y:A$ 以及 $n\ge -2$，映射~\eqref{eq:path-trunc-map} 是一个等价；因此我们有
  \[ \eqv{\ttrunc n{x=_Ay}}{\Big(\tproj {n+1}x=_{\trunc{n+1}A}\tproj {n+1}y\Big)}. \]
\end{thm}

\begin{proof}
  证明是编码-解码方法的一个简单应用：
  如同以往的情形一样，我们无法直接定义映射~\eqref{eq:path-trunc-map} 的一个拟逆，因为我们无法对 $\tproj {n+1}x$ 和 $\tproj {n+1}y$ 之间的相等进行归纳。
  因此，我们转而将其类型推广，以便处理类型 $\trunc{n+1}A$ 中的一般元素，而不仅仅是 $\tproj {n+1}x$ 和 $\tproj {n+1}y$。
  如下定义 $P:\trunc {n+1}A\to\trunc {n+1}A\to\typele{n}$：
  \[P(\tproj {n+1}x,\tproj {n+1}y)\defeq \trunc n{x=_Ay}\]
  这个定义是良定义的，因为 $\trunc n{x=_Ay}$ 是 $n$-截断的，而根据 \cref{thm:hleveln-of-hlevelSn}，$\typele{n}$ 是 $(n+1)$-截断的。
  现在，对于每一对 $u,v:\trunc{n+1}A$，都有一个映射
  \[\decode:P(u,v) \to \big(u=_{\trunc{n+1}A}v\big)\]
  它对 $u=\tproj {n+1}x$、$v=\tproj {n+1}y$ 以及 $p:x=y$ 定义为
  \[\decode(\tproj n{p})\defeq \apfunc{\tprojf{n+1}} (p).\]
  由于 $\decode$ 的余域是 $n$-截断的，我们只需对 $u$ 和 $v$ 是这种形式的情形定义它即可，而这与之前的定义相同。
  我们还通过对 $u$ 作归纳来定义一个函数
  \[ r : \prd{u:\trunc{n+1} A} P(u,u) \]
  其中 $r(\tproj{n+1} x) \defeq \tproj n {\refl x}$。

  现在我们可以定义一个逆映射
  \[\encode: (u=_{\trunc{n+1}A}v) \to P(u,v)\]
  定义为
  \[\encode(p) \defeq \transfib{v\mapsto P(u,v)}{p}{r(u)}. \]
  为了证明复合映射
  \[ (u=_{\trunc{n+1}A}v) \xrightarrow{\encode} P(u,v) \xrightarrow{\decode} (u=_{\trunc{n+1}A}v) \]
  是恒等函数，根据道路归纳，只需对 $\refl u : u=u$ 的情形进行验证；此时我们需要知道的是 $\decode(r(u)) = \refl{u}$。
  但由于这是一个 $(n-1)$-类型，因而也是一个 $(n+1)$-类型，我们可以假定 $u\jdeq \tproj {n+1} x$，而在此假定下，结论由 $r$ 和 $\decode$ 的定义即得。
  最后，为了证明
  \[ P(u,v) \xrightarrow{\decode} (u=_{\trunc{n+1}A}v) \xrightarrow{\encode} P(u,v) \]
  是恒等函数，由于该目标同样是一个 $(n-1)$-类型，我们可以假定 $u=\tproj {n+1}x$ 且 $v=\tproj {n+1}y$，且我们考虑的是对应于某个 $p:x=y$ 的 $\tproj n p:P(\tproj{n+1}x,\tproj{n+1}y)$。
  那么我们有
  \begin{align*}
    \encode(\decode(\tproj n p)) &= \encode(\apfunc{\tprojf{n+1}}(p))\\
    &= \transfib{v\mapsto P(\tproj{n+1}x,v)}{\apfunc{\tprojf{n+1}}(p)}{\tproj n {\refl x}}\\
    &= \transfib{y\mapsto \trunc n{x=y}}{p}{\tproj n {\refl x}}\\
    &= \tproj n {\transfib{y \mapsto (x=y)}{p}{\refl x}}\\
    &= \tproj n p,
  \end{align*}
  这里用到了 \cref{thm:transport-compose,thm:ap-transport}。
  （或者，我们也可以对 $p$ 进行道路归纳；则所需的相等在判断意义下成立。）
  这就完成了 \decode 和 \encode 是拟逆的证明。
  所陈述的结果便是 $u=\tproj {n+1}x$ 且 $v=\tproj {n+1}y$ 时的特例。
\end{proof}

\begin{cor}
  设 $n\ge-2$，且 $(A,a)$ 是一个带点类型。那么
  \[\ttrunc n{\Omega(A,a)}=\Omega\mathopen{}\left(\trunc{n+1}{(A,a)}\right)\]
\end{cor}

\begin{proof}
  这是前一个引理在 $x=y=a$ 时的特例。
\end{proof}

\begin{cor}
  设 $n\ge -2$，$k\ge 0$，且 $(A,a)$ 是一个带点类型。那么
  \[\ttrunc n{\Omega^k(A,a)} = \Omega^k\mathopen{}\left(\trunc{n+k}{(A,a)}\right). \]
\end{cor}

\begin{proof}
  对 $k$ 进行归纳，利用 $\Omega^k$ 的递归定义即可。
\end{proof}

我们还注意到“截断是累积的”：如果我们将一个类型先截断为一个 $n$-类型，再截断为一个 $k$-类型（其中 $k\le n$），那么我们不妨一开始就直接将其截断为一个 $k$-类型。

\begin{lem} \label{lem:truncation-le}
  设 $k,n\ge-2$ 且 $k\le{}n$，$A:\type$。那么
  $\trunc k{\trunc nA}=\trunc kA$。
\end{lem}

\begin{proof}
  我们定义两个映射 $f:\trunc k{\trunc nA}\to\trunc kA$ 和
  $g:\trunc kA\to\trunc k{\trunc nA}$ 如下：
  %
  \[
   f(\tproj k{\tproj na}) \defeq \tproj ka
   \qquad\text{and}\qquad
   g(\tproj ka) \defeq \tproj k{\tproj na}.
  \]
  %
  映射 $f$ 是良定义的，因为 $\trunc kA$ 是 $k$-截断的，并且也是 $n$-截断的（因为 $k\le{}n$）；映射 $g$ 是良定义的，因为 $\trunc k{\trunc nA}$ 是 $k$-截断的。

  复合映射 $f\circ{}g:\trunc kA\to\trunc kA$ 满足
  $(f\circ{}g)(\tproj ka)=\tproj ka$，因此 $f\circ{}g=\idfunc[\trunc kA]$。
  类似地，我们有 $(g\circ{}f)(\tproj k{\tproj na})=\tproj k{\tproj na}$，因此 $g\circ{}f=\idfunc[\trunc k{\trunc nA}]$。
\end{proof}

% \begin{lem}
%   We have $\trunc n{\unit}=\unit$.
% \end{lem}
% \begin{proof}
%   Indeed, $\unit$ is $n$-truncated for every $n$ hence $\trunc n{\unit}=\unit$ by
%   \cref{reflectPequiv}.
% \end{proof}

\index{truncation!n-truncation@$n$-truncation|)}%

\section{$n$-类型的余极限}
\label{sec:pushouts}

回想一下，在 \cref{sec:colimits} 中，我们曾用高阶归纳类型定义了类型的推出，并证明了其泛性质。
一般来说，$n$-\type{} 的（同伦）余极限可能不再是 $n$-\type{}（关于一个极端的反例，可见 \cref{ex:s2-colim-unit}）。
然而，如果对其进行 $n$-截断，我们便能得到一个 $n$-\type{}，并且它对于其他 $n$-\type{} 而言满足正确的泛性质。

本节我们将对推出这一情形证明上述结论，因为推出是余极限中最重要且非平凡的情形。
我们先回顾 \cref{sec:colimits} 中给出的如下定义。

\begin{defn}
  一个 \define{伸展（span）}% in $\P$
  \indexdef{span}%
  由一个五元组 $\Ddiag=(A,B,C,f,g)$ 构成，其中 % $A,B,C:\P$ and，并有映射 $f:C\to{}A$ 与 $g:C\to{}B$。
  \[\Ddiag=\quad\vcenter{\xymatrix{C \ar^g[r] \ar_f[d] & B \\ A & }}\]
\end{defn}

\begin{defn}
  给定一个伸展 $\Ddiag=(A,B,C,f,g)$ 和一个类型 $D$，%$D:\P$, a
  \define{上锥（cocone under $\Ddiag$ with base $D$）} 是一个三元组 $(i, j, h)$，
  \index{cocone}%
  其中 $i:A\to{}D$， $j:B\to{}D$，并且 $h : \prd{c:C}i(f(c))=j(g(c))$：
  \[\uppercurveobject{{ }}\lowercurveobject{{ }}\twocellhead{{ }}
  \xymatrix{C \ar^g[r] \ar_f[d] \drtwocell{^h} & B \ar^j[d] \\ A \ar_i[r] & D
  }\]
  我们用 $\cocone{\Ddiag}{D}$ 表示所有这类上锥构成的类型。
\end{defn}

上锥的类型具有（协变）函子性。
例如，给定 $D$， $E$% $D,E:\P$
以及一个映射 $t:D\to{}E$，则存在一个映射
  \[\function{\cocone{\Ddiag}{D}}{\cocone{\Ddiag}{E}}{c}{\composecocone{t}c}\]
  其定义为：
  \[\composecocone{t}(i,j,h)=(t\circ{}i,t\circ{}j,\mapfunc{t}\circ{}h).\]
而给定 $D$， $E$， $F$，%$:\P$,
函数 $t:D\to{}E$， $u:E\to{}F$ 以及 $c:\cocone{\Ddiag}{D}$，则有
\begin{align}
  \composecocone{\idfunc[D]}c &= c \label{eq:composeconeid}\\
  \composecocone{(u\circ{}t)}c&=\composecocone{u}(\composecocone{t}c). \label{eq:composeconefunc}
\end{align}

\begin{defn}
  给定一个由 $n$-\type{} 构成的伸展 $\Ddiag$，一个 $n$-\type{} $D$，以及一个上锥
  $c:\cocone{\Ddiag}{D}$，如果对每个 $n$-\type{} $E$，映射
  \[\function{(D\to{}E)}{\cocone{\Ddiag}{E}}{t}{\composecocone{t}c}\]
  都是一个等价，则称配对 $(D,c)$ 为 \define{在 $n$-\type{} 中的推出（pushout of $\Ddiag$ in $n$-types）}。
  \indexdef{pushout!in ntypes@in $n$-types}%
\end{defn}

\begin{comment}
We showed in \cref{thm:pushout-ump} that pushouts exist when $\P$ is \type itself, by giving a direct construction in terms of higher
inductive types.
For a general \P, pushouts may or may not exist, but if they do, then they are unique.

\begin{lem}
  If $(D,c)$ and $(D',c')$ are two pushouts of $\Ddiag$ in $\P$, then
  $(D,c)=(D',c')$.
\end{lem}
\begin{proof}
  We first prove that the two types $D$ and $D'$ are equivalent.

  Using the universal property of $D$ with $D'$, we see that the following map is an
  equivalence
  %
  \[
    \function{(D\to{}D')}{\cocone{\Ddiag}{D'}}{t}{\composecocone{t}c}
  \]
  %
  In particular, there is a function $f:D\to{}D'$ satisfying $\composecocone{f}c=c'$. In the
  same way there is a function $g:D'\to{}D$ such that $\composecocone{g}c'=c$.

  In order to prove that $g\circ{}f=\idfunc[D]$ we use the universal property of
  $D$ for $D$, which says that the following map is an equivalence:
  %
  \[
  \function{(D\to{}D)}{\cocone{\Ddiag}{D}}{t}{\composecocone{t}c}
  \]
  %
  Using the functoriality of $t\mapsto{}\composecocone{t}c$ we see that
  \begin{align*}
    \composecocone{(g\circ{}f)}c &= \composecocone{g}(\composecocone{f}c) \\
    &= \composecocone{g}c' \\
    &= c \\
    &= \composecocone{\idfunc[D]}c
  \end{align*}
  hence
  $g\circ{}f=\idfunc[D]$, because equivalences are injective. The same argument
  with $D'$ instead of $D$ shows that $f\circ{}g=\idfunc[D']$.

  Hence $D$ and $D'$ are equal, and the fact that $(D,c)=(D',c')$ follows from
  the fact that the equivalence between $D$ and $D'$ we just defined sends $c$
  to $c'$.
\end{proof}

\begin{cor}
  The type of pushouts of $\Ddiag$ in $\P$ is a mere proposition. In particular if
  pushouts merely exist then they actually exist.
\end{cor}

As in the case of pullbacks, if \P is reflective, then pushouts in \P always exist.
However, unlike the case of pullbacks, pushouts in \P are not the same as the pushouts in \type: they are obtained by applying the
reflector.
\end{comment}

为了构造 $n$-\type{} 的推出，我们需要先说明如何反射伸展与上锥。

\bgroup
\def\reflect(#1){\trunc n{#1}}

\begin{defn}
  设
  \[\Ddiag=\quad\vcenter{\xymatrix{C \ar^g[r] \ar_f[d] & B \\ A & }}\]
  为一个伸展（span）。我们用 $\reflect(\Ddiag)$ 表示如下的 $n$-类型伸展：
  \[\reflect(\Ddiag)\defeq\quad \vcenter{\xymatrix{\reflect(C) \ar^{\reflect(g)}[r]
      \ar_{\reflect(f)}[d] & \reflect(B) \\ \reflect(A) & }}\]
\end{defn}

\begin{defn}
  设 $D:\type$ 且 $c=(i,j,h):\cocone{\Ddiag}{D}$。
  我们定义
  \[\reflect(c)=(\reflect(i),\reflect(j),k):
  \cocone{\reflect(\Ddiag)}{\reflect(D)}\]
  其中 $k$ 为复合同伦
  \[ \reflect(i) \circ \reflect(f) \htpy \reflect(i\circ f) \htpy \reflect(j\circ g) \htpy \reflect(j) \circ \reflect(g) \]
  这里使用了 \cref{thm:trunc-htpy} 以及 $\reflect(\blank)$ 的函子性。
  % \[\reflect(h):\prd{c:\reflect(C)}\reflect(i)(\reflect(f)(c))=\reflect(j)(\reflect(g)(c))\]
  % is defined in the following way:
\end{defn}

\egroup

我们现在观察到，从每个类型到其 $n$-截断的映射，在以下意义下组合成了一个伸展之间的映射。

\begin{defn}
  设
  \[\Ddiag=\quad\vcenter{\xymatrix{C \ar^g[r] \ar_f[d] & B \\ A & }}
  \qquad\text{and}\qquad
  \Ddiag'=\quad\vcenter{\xymatrix{C' \ar^{g'}[r] \ar_{f'}[d] & B' \\ A' & }}
  \]
  为两个伸展。
  一个从 $\Ddiag$ 到 $\Ddiag'$ 的\define{伸展之间的映射}
  \indexdef{map!of spans}%
  由函数 $\alpha:A\to A'$、$\beta:B\to B'$、$\gamma:C\to C'$ 以及同伦 $\phi: \alpha\circ f \htpy f'\circ \gamma$ 和 $\psi:\beta\circ g \htpy g' \circ \gamma$ 构成。
\end{defn}

因此，对任意伸展 $\Ddiag$，我们都有一个由 $\tprojf[A]n$、$\tprojf[B]n$、$\tprojf[C]n$ 以及来自~\eqref{eq:trunc-nat} 的自然性同伦 $\mathsf{nat}^f_n$ 和 $\mathsf{nat}^g_n$ 所构成的伸展之间的映射 $\tprojf[\Ddiag] n : \Ddiag \to \trunc n\Ddiag$。

我们还需要知道伸展之间的映射具有函子性。
即，若 $(\alpha,\beta,\gamma,\phi,\psi):\Ddiag \to \Ddiag'$ 是一个伸展之间的映射且 $D$ 是任意类型，则有
\[ \function{\cocone{\Ddiag'}{D}}{\cocone{\Ddiag}{D}}{(i,j,h)}{(i\circ \alpha,j\circ\beta, k)} \]
其中 $k: \prd{z:C} i(\alpha(f(z))) = j(\beta(g(z)))$ 是复合
\begin{equation}\label{eq:mapofspans-htpy}
\xymatrix{
  i(\alpha(f(z))) \ar@{=}[r]^{\apfunc{i}(\phi)} &
  i(f'(\gamma(z))) \ar@{=}[r]^{h(\gamma(z))} &
  j(g'(\gamma(z))) \ar@{=}[r]^{\apfunc{j}(\psi)} &
  j(\beta(g(z))). }
\end{equation}
我们将这个上锥记为 $(i,j,h) \circ (\alpha,\beta,\gamma,\phi,\psi)$。
此外，该函子作用与上锥的另一函子作用可交换：

\begin{lem}\label{thm:conemap-funct}
  给定 $(\alpha,\beta,\gamma,\phi,\psi):\Ddiag \to \Ddiag'$ 以及 $t:D\to E$，下图交换：
  \begin{equation*}
    \vcenter{\xymatrix{
        \cocone{\Ddiag'}{D}\ar[r]^{t \circ {\blank}}\ar[d] &
        \cocone{\Ddiag'}{E}\ar[d]\\
        \cocone{\Ddiag}{D}\ar[r]_{t \circ {\blank}} &
        \cocone{\Ddiag}{E}
      }}
  \end{equation*}
\end{lem}

\begin{proof}
  给定 $(i,j,h):\cocone{\Ddiag'}{D}$，注意两个复合都给出一个上锥，其前两个分量为 $t\circ i\circ \alpha$ 和 $t\circ j\circ\beta$。
  因此，只需验证二者的同伦是否相同。
  对于右上方的复合，其同伦是将 $(i,j,h)$ 替换为 $(t\circ i, t\circ j, \apfunc{t}\circ h)$ 后的~\eqref{eq:mapofspans-htpy}：
  \begin{equation*}
    \xymatrix@+2.8em{
      {t \, i \, \alpha \, f \, z} \ar@{=}[r]^{\apfunc{t\circ i}(\phi)} &
      {t \, i \, f' \, \gamma \, z} \ar@{=}[r]^{\apfunc{t}(h(\gamma(z)))} &
      {t \, j \, g' \, \gamma \, z} \ar@{=}[r]^{\apfunc{t\circ j}(\psi)} &
      {t \, j \, \beta \, g \, z}
    }
  \end{equation*}
  （为简洁起见，我们省略了函数参数周围的括号。）
  另一方面，对于左下方的复合，其同伦是 $\apfunc{t}$ 作用于~\eqref{eq:mapofspans-htpy} 的结果。
  由于 $\apfunc{}$ 保路径的拼接，这等于
  \begin{equation*}
    \xymatrix@+2.8em{
      {t \, i \, \alpha \, f \, z} \ar@{=}[r]^{\apfunc{t}(\apfunc{i}(\phi))} &
      {t \, i \, f' \, \gamma \, z} \ar@{=}[r]^{\apfunc{t}(h(\gamma(z)))} &
      {t \, j \, g' \, \gamma \, z} \ar@{=}[r]^{\apfunc{t}(\apfunc{j}(\psi))} &
      {t \, j \, \beta \, g \, z}. }
  \end{equation*}
  但是 $\apfunc{t}\circ \apfunc{i} = \apfunc{t\circ i}$，且对 $j$ 类似，因此这两个同伦相等。
\end{proof}

最后注意到，由于我们使用 \cref{thm:trunc-htpy} 定义了 $\trunc nc : \cocone{\trunc n \Ddiag}{\trunc n D}$，附加条件~\eqref{eq:trunc-htpy} 蕴含着
\begin{equation}
  \tprojf[D] n \circ c = \trunc n c \circ \tprojf[\Ddiag]n. \label{eq:conetrunc}
\end{equation}
对任意 $c:\cocone{\Ddiag}{D}$ 成立。
现在我们可以证明期望的定理了。

\begin{thm}
  \label{reflectcommutespushout}
  \index{universal!property!of pushout}%
  设 $\Ddiag$ 为一个伸展，$(D,c)$ 是其推出。
  则 $(\trunc nD,\trunc n c)$ 是 $\trunc n\Ddiag$ 在 $n$-类型中的推出。
\end{thm}

\begin{proof}
  设 $E$ 是一个 $n$-类型，考虑下图：
\bgroup
\def\reflect(#1){\trunc n{#1}}
  \begin{equation*}
  \vcenter{\xymatrix{
      (\trunc nD \to E)\ar[r]^-{\blank\circ \tprojf[D] n}\ar[d]_{\blank\circ \trunc nc} &
      (D\to E)\ar[d]^{\blank\circ c}\\
      \cocone{\trunc n \Ddiag}{E}\ar[r]^-{\blank\circ \tprojf[\Ddiag]n}\ar@{<-}[d]_{\ell_1} &
      \cocone{\Ddiag}{E}\ar@{<-}[d]^{\ell_2}\\
      (\reflect(A)\to{}E)\times_{(\reflect(C)\to{}E)}(\reflect(B)\to{}E)\ar[r] &
      (A\to{}E)\times_{(C\to{}E)}(B\to{}E)
      }}
  \end{equation*}
\egroup
  上方水平箭头是等价，因为 $E$ 是 $n$-类型；而 $\blank\circ c$ 是等价，因为 $c$ 是推出上锥。
  因此，根据 2-out-of-3 性质，要证明 $\blank\circ \trunc nc$ 是等价，只需证明上方方块交换且中间水平箭头是等价。
  为证明上方方块交换，令 $t:\trunc nD \to E$；则有
  \begin{align}
    \big(t \circ \trunc n c\big) \circ \tprojf[\Ddiag] n
    &= t \circ \big(\trunc n c \circ \tprojf[\Ddiag] n\big)
    \tag{by \cref{thm:conemap-funct}}\\
    &= t\circ \big(\tprojf[D]n \circ c\big)
    \tag{by~\eqref{eq:conetrunc}}\\
    &= \big(t\circ \tprojf[D]n\big) \circ c
    \tag{by~\eqref{eq:composeconefunc}}.
  \end{align}
  为证明中间水平箭头是等价，考虑下方方块。
  下方两个垂直箭头就是 $\happly$ 的应用：
  \begin{align*}
    \ell_1(i,j,p) &\defeq (i,j,\happly(p))\\
    \ell_2(i,j,p) &\defeq (i,j,\happly(p))
  \end{align*}
  因此由函数外延性可知它们是等价。
  最下方的水平箭头定义为
  \[ (i,j,p) \mapsto \big( i\circ \tprojf[A]n,\;\; j \circ \tprojf[B] n,\;\; q\big) \]
  其中 $q$ 是复合
  \begin{align}
    i\circ \tprojf[A]n \circ f
    &= i\circ \trunc nf \circ \tprojf[C]n
    \tag{by $\funext(\lam{z} \apfunc{i}(\mathsf{nat}^f_n(z)))$}\\
    &= j\circ \trunc ng \circ \tprojf[C]n
    \tag{by $\apfunc{\blank\circ \tprojf[C] n}(p)$}\\
    &= j\circ \tprojf[B]n \circ g.
    \tag{by $\funext(\lam{z} \apfunc{j}(\mathsf{nat}^g_n(z)))$}
  \end{align}
  这是一个等价，因为它由一个余伸展的等价所诱导。
  因此，再次由 2-out-of-3，只需证明下方方块交换。
  但绕行下方方块的两个复合在前两个分量上依定义相等，故只需证明对左下角的 $(i,j,p)$ 和 $z:C$，路径
  \[ \happly(q,z) : i(\tproj n{f(z)}) = j(\tproj n{g(z)}) \]
  （其中 $q$ 如上所述）
  等于复合
  \begin{align}
    i(\tproj n{f(z)})
    &= i(\trunc nf(\tproj nz))
    \tag{by $\apfunc{i}(\mathsf{nat}^f_n(z))$}\\
    &= j(\trunc ng(\tproj nz))
    \tag{by $\happly(p,\tproj nz)$}\\
    &= j(\tproj n{g(z)}).
    \tag{by $\apfunc{j}(\mathsf{nat}^g_n(z))$}
  \end{align}
  然而，由于 $\happly$ 具有函子性，只需检查三个分量路径相等即可：
  \begin{align*}
    \happly({\funext(\lam{z} \apfunc{i}(\mathsf{nat}^f_n(z)))},z)
    &= {\apfunc{i}(\mathsf{nat}^f_n(z))}\\
    \happly(\apfunc{\blank\circ \tprojf[C] n}(p), z)
    &= {\happly(p,\tproj nz)}\\
    \happly({\funext(\lam{z} \apfunc{j}(\mathsf{nat}^g_n(z)))},z)
    &= {\apfunc{j}(\mathsf{nat}^g_n(z))}.
  \end{align*}
  其中第一和第三条路径正是 $\happly$ 是 $\funext$ 的拟逆这一事实，而第二条路径则是关于 $\happly$ 与预复合的一个简单通用引理。
\end{proof}

\section{连通性}
\label{sec:connectivity}

$n$-类型是指在 $n$ 维以上不包含有趣信息的类型。
与之相对，一个\emph{$n$-连通类型}是指在 $n$ 维\emph{以下}不包含有趣信息的类型。
事实证明，研究函数的更一般概念也是很自然的。

\begin{defn}
一个函数 $f:A\to B$ 被称为是 \define{$n$-连通的（$n$-connected）}
\indexdef{function!n-connected@$n$-connected}%
\indexsee{n-connected@$n$-connected!function}{function, $n$-connected}%
如果对于所有 $b:B$，类型 $\trunc n{\hfiber f b}$ 都是可缩的：
\begin{equation*}
  \mathsf{conn}_n(f)\defeq \prd{b:B}\iscontr(\trunc n{\hfiber{f}b}).
\end{equation*}
一个类型 $A$ 被称为是 \define{$n$-连通的（$n$-connected）}
\indexsee{n-connected@$n$-connected!type}{type, $n$-connected}%
\indexdef{type!n-connected@$n$-connected}%
 如果唯一的函数 $A\to\unit$ 是 $n$-连通的，亦即，如果 $\trunc nA$ 是可缩的。
\end{defn}


\indexsee{connected!function}{function, $n$-connected}

因此，一个函数 $f:A\to B$ 是 $n$-连通的，当且仅当对每个 $b:B$，其纤维 $\hfib{f}b$ 都是 $n$-连通的。
显然，每个函数都是 $(-2)$-连通的。
在下一个层级上，我们有：

\begin{lem}\label{thm:minusoneconn-surjective}
  \index{function!surjective}%
  一个函数 $f$ 是 $(-1)$-连通的，当且仅当它是\cref{sec:mono-surj}意义下的满的（surjective）。
\end{lem}

\begin{proof}
  我们曾定义 $f$ 是满的，若对所有 $b$，类型 $\trunc{-1}{\hfiber f b}$ 都有实例。
  但由于这是一个命题，有实例等价于可缩。
\end{proof}

这样一来，对于 $n\ge 0$，函数的 $n$-连通性可以看作是一种更强的满射性。
从范畴论角度看，$(-1)$-连通性对应于对象上的本质满性，而 $n$-连通性则对应于 $k$-态射（$k\le n+1$）上的本质满性。

\cref{thm:minusoneconn-surjective}还蕴含着，一个类型 $A$ 是 $(-1)$-连通的，当且仅当它仅知有实例。
当一个类型是 $0$-连通的时候，我们可以简称为它是 \define{连通的}，
\indexdef{connected!type}%
\indexdef{type!connected}%
而当它是 $1$-连通的，我们则称之为是 \define{单连通的}。
\indexdef{simply connected type}%
\indexdef{type!simply connected}%

\begin{rmk}\label{rmk:connectedness-indexing}
  尽管我们对于类型的 $n$-连通性定义与同伦论中的标准定义一致，但我们对于\emph{函数}的 $n$-连通性定义，与经典同伦论中一个常见的编号约定相差了 1。
  我们称一个函数 $f$ 是 $n$-连通的，如果它的所有纤维都是 $n$-连通的；而一些经典同伦论学者会称这样的函数为 $(n+1)$-连通的。
  （这是由于历史上更关注\emph{余纤维（cofibers）}而非纤维。）
\end{rmk}

现在我们来观察连通映射的几条闭包性质。
\index{function!n-connected@$n$-connected}

\begin{lem}
\index{retract!of a function}%
设 $g$ 是一个 $n$-连通函数 $f$ 的收缩核。那么 $g$ 也是 $n$-连通的。
\end{lem}

\begin{proof}
这是\cref{lem:func_retract_to_fiber_retract}的直接推论。
\end{proof}

\begin{cor}
如果 $g$ 同伦于一个 $n$-连通函数 $f$，那么 $g$ 也是 $n$-连通的。
\end{cor}

\begin{lem}\label{lem:nconnected_postcomp}
设 $f:A\to B$ 是 $n$-连通的。那么 $g:B\to C$ 是 $n$-连通的，当且仅当复合 $g\circ f$ 是 $n$-连通的。
\end{lem}

\begin{proof}
对任意 $c:C$，我们有
\begin{align*}
  \trunc n{\hfib{g\circ f}c}
  & \eqvsym \Trunc n{ \sm{w:\hfib{g}c}\hfib{f}{\proj1 w}}
  \tag{by \cref{ex:unstable-octahedron}}\\
  & \eqvsym \Trunc n{\sm{w:\hfib{g}c} \trunc n{\hfib{f}{\proj1 w}}}
  \tag{by \cref{thm:trunc-in-truncated-sigma}}\\
  & \eqvsym \trunc n{\hfib{g}c}.
  \tag{since $\trunc n{\hfib{f}{\proj1 w}}$ is contractible}
\end{align*}
由此可见，$\trunc n{\hfib{g}c}$ 是可缩的，当且仅当$\trunc n{\hfib{g\circ f}c}$是可缩的。
\end{proof}

重要的是，$n$-连通函数可以等价地刻画为那些满足关于 $n$-类型的“归纳原理”的函数。\index{induction principle!for connected maps}
这一想法将直接引导我们在\cref{sec:freudenthal}中证明 Freudenthal 纬悬定理。

\begin{lem}\label{prop:nconnected_tested_by_lv_n_dependent types}
对于函数 $f:A\to B$ 和类型族 $P:B\to\type$，考虑以下函数：
\begin{equation*}
\lam{s} s\circ f :\Parens{\prd{b:B} P(b)}\to\Parens{\prd{a:A}P(f(a))}.
\end{equation*}
对于一个固定的 $f$ 和 $n\ge -2$，下列命题等价：
\begin{enumerate}
\item $f$ 是 $n$-连通的。\label{item:conntest1}
\item 对于每个 $P:B\to\ntype{n}$，映射 $\lam{s} s\circ f$ 都是一个等价。\label{item:conntest2}
\item 对于每个 $P:B\to\ntype{n}$，映射 $\lam{s} s\circ f$ 都有一个截面。\label{item:conntest3}
\end{enumerate}
\end{lem}

\begin{proof}
假设 $f$ 是 $n$-连通的，并设 $P:B\to\ntype{n}$。那么我们有如下等价关系：
\begin{align}
  \prd{b:B} P(b) & \eqvsym \prd{b:B} \Parens{\trunc n{\hfib{f}b} \to P(b)}
  \tag{since $\trunc n{\hfib{f}b}$ is contractible}\\
  & \eqvsym \prd{b:B} \Parens{\hfib{f}b\to P(b)}
  \tag{since $P(b)$ is an $n$-type}\\
  & \eqvsym \prd{b:B}{a:A}{p:f(a)= b} P(b)
  \tag{by the left universal property of $\Sigma$-types}\\
  & \eqvsym \prd{a:A} P(f(a)).
  \tag{by the left universal property of path types}
\end{align}
我们略去关于此等价确实由 $\lam{s} s\circ f$ 给出的证明。
因此，有\ref{item:conntest1}$\Rightarrow$\ref{item:conntest2}，且显然有\ref{item:conntest2}$\Rightarrow$\ref{item:conntest3}。
为证\ref{item:conntest3}$\Rightarrow$\ref{item:conntest1}，考虑类型族
\begin{equation*}
P(b)\defeq \trunc n{\hfib{f}b}.
\end{equation*}
那么，由\ref{item:conntest3}可得到一个映射$c:\prd{b:B} \trunc n{\hfib{f}b}$，满足$c(f(a))=\tproj n{\pairr{a,\refl{f(a)}}}$。
为证明每个 $\trunc n{\hfib{f}b}$ 都是可缩的，我们需要构造一个类型如下的函数：
\begin{equation*}
\prd{b:B}{w:\trunc n{\hfib{f}b}} w= c(b).
\end{equation*}
根据\cref{thm:truncn-ind}，为此只需构造一个类型如下的函数：
\begin{equation*}
\prd{b:B}{a:A}{p:f(a)= b} \tproj n{\pairr{a,p}}= c(b).
\end{equation*}
但通过调整变量并使用道路归纳，这等价于如下类型：
\begin{equation*}
\prd{a:A} \tproj n{\pairr{a,\refl{f(a)}}}= c(f(a)).
\end{equation*}
而根据我们对 $c(f(a))$ 的选取，这一性质成立。
\end{proof}

\begin{cor}\label{cor:totrunc-is-connected}
对任意类型 $A$，典范函数 $\tprojf n:A\to\trunc n A$ 是 $n$-连通的。
\end{cor}

\begin{proof}
根据\cref{thm:truncn-ind}以及与之相伴的唯一性原理，\cref{prop:nconnected_tested_by_lv_n_dependent types}的条件成立。
\end{proof}

例如，当 $n=-1$ 时，\cref{cor:totrunc-is-connected}说明从一个类型到其命题截断的映射 $A\to \brck A$ 是满的。

\begin{cor}\label{thm:nconn-to-ntype-const}\label{connectedtotruncated}
一个类型 $A$ 是 $n$-连通的，当且仅当对每个 $n$-类型 $B$，映射
\begin{equation*}
  \lam{b}{a} b: B \to (A\to B)
\end{equation*}
都是等价。
换言之，“从 $A$ 到任何 $n$-类型的映射都是常值映射”。
\end{cor}

\begin{proof}
  将 \cref{prop:nconnected_tested_by_lv_n_dependent types} 应用于一个陪域为 $\unit$ 的函数即得。
\end{proof}

\begin{lem}\label{lem:nconnected_to_leveln_to_equiv}
设 $B$ 是一个 $n$-类型，并设 $f:A\to B$ 是一个函数。那么诱导函数 $g:\trunc n A\to B$ 是等价，当且仅当 $f$ 是 $n$-连通的。
\end{lem}

\begin{proof}
根据 \cref{cor:totrunc-is-connected}，$\tprojf n$ 是 $n$-连通的。
由于 $f = g\circ \tprojf n$，根据
\cref{lem:nconnected_postcomp} 可知，$f$ 是 $n$-连通的当且仅当 $g$ 是 $n$-连通的。
但 $g$ 是两个 $n$-类型之间的函数，所以它的纤维也都是 $n$-类型。
因此，$g$ 是 $n$-连通的当且仅当它是等价。
\end{proof}

我们也可以用基点包含映射的连通性来刻画连通的带点类型。

\begin{lem}\label{thm:connected-pointed}
  \index{basepoint}%
  设 $A$ 是一个类型，$a_0:\unit\to A$ 是一个基点，且 $n\ge -1$。
  那么 $A$ 是 $n$-连通的，当且仅当映射 $a_0$ 是 $(n-1)$-连通的。
\end{lem}

\begin{proof}
  首先假设 $a_0:\unit\to A$ 是 $(n-1)$-连通的，并设 $B$ 是一个 $n$-类型；我们将使用 \cref{thm:nconn-to-ntype-const}。
  映射 $\lam{b}{a} b: B \to (A\to B)$ 有一个收缩，由 $f\mapsto f(a_0)$ 给出，因此只需证明它也有一个截面，即对任意 $f:A\to B$，存在 $b:B$ 使得 $f = \lam{a}b$。
  我们取 $b\defeq f(a_0)$。
  定义 $P:A\to\type$ 为 $P(a) \defeq (f(a)=f(a_0))$。
  那么 $P$ 是一个由 $(n-1)$-类型构成的族，并且我们有 $P(a_0)$；由于 $a_0:\unit\to A$ 是 $(n-1)$-连通的，因此我们有 $\prd{a:A} P(a)$。
  于是，$f = \lam{a} f(a_0)$ 如所需。

  现在假设 $A$ 是 $n$-连通的，设给定 $P:A\to\ntype{(n-1)}$ 和 $u:P(a_0)$。
  根据 \cref{prop:nconnected_tested_by_lv_n_dependent types}，只需构造 $f:\prd{a:A} P(a)$ 使得 $f(a_0)=u$。
  现在 $\ntype{(n-1)}$ 是一个 $n$-类型且 $A$ 是 $n$-连通的，因此由 \cref{thm:nconn-to-ntype-const}，存在一个 $n$-类型 $B$ 使得 $P = \lam{a} B$。
  于是，我们得到一族等价 $g:\prd{a:A} (\eqv{P(a)}{B})$。
  定义 $f(a) \defeq \opp{g_a}(g_{a_0}(u))$；则 $f:\prd{a:A} P(a)$ 且 $f(a_0) = u$ 如所需。
\end{proof}

特别地，一个带点类型 $(A,a_0)$ 是 0-连通的，当且仅当 $a_0:\unit\to A$ 是满的，也就是说 $\prd{x:A} \brck{x=a_0}$。
关于不一定带点的情形的类似结果，参见 \cref{ex:connectivity-inductively}。

\cref{lem:nconnected_postcomp} 的一个有用变体是：

\begin{lem}\label{lem:nconnected_postcomp_variation}
设 $f:A\to B$ 是一个函数，$P:A\to\type$ 和 $Q:B\to\type$ 是类型族。假设 $g:\prd{a:A} P(a)\to Q(f(a))$
是一个逐纤维 $n$-连通%
\index{fiberwise!n-connected family of functions@$n$-connected family of functions}
的函数族，即每个函数 $g_a : P(a) \to Q(f(a))$ 都是 $n$-连通的。
若 $f$ 也是 $n$-连通的，则函数
\begin{align*}
\varphi &:\Parens{\sm{a:A} P(a)}\to\Parens{\sm{b:B} Q(b)}\\
\varphi(a,u) &\defeq \pairr{f(a),g_a(u)}.
\end{align*}
也是 $n$-连通的。
反之，若 $\varphi$ 和每个 $g_a$ 都是 $n$-连通的，且 $Q$ 逐纤维地仅知有实例（即对所有 $b:B$ 有 $\brck{Q(b)}$），则 $f$ 是 $n$-连通的。
\end{lem}

\begin{proof}
对任意 $b:B$ 和 $v:Q(b)$，我们有
{\allowdisplaybreaks
\begin{align*}
\trunc n{\hfib{\varphi}{\pairr{b,v}}} & \eqvsym \Trunc n{\sm{a:A}{u:P(a)}{p:f(a)= b} \trans{p}{g_a(u)}= v}\\
& \eqvsym \Trunc n{\sm{w:\hfib{f}b}{u:P(\proj1(w))} g_{\proj 1 w}(u)= \trans{\opp{\proj2(w)}}{v}}\\
& \eqvsym \Trunc n{\sm{w:\hfib{f}b} \hfib{g(\proj1 w)}{\trans{\opp{\proj 2(w)}}{v}}}\\
& \eqvsym \Trunc n{\sm{w:\hfib{f}b} \trunc n{\hfib{g(\proj1 w)}{\trans{\opp{\proj 2(w)}}{v}}}}\\
& \eqvsym \trunc n{\hfib{f}b}
\end{align*}}%
其中沿 $f(p)$ 和 $f(p)^{-1}$ 的输运是关于 $Q$ 的。
因此，若其中一个是可缩的，则另一个也是可缩的。

特别地，若 $f$ 是 $n$-连通的，则 $\trunc n{\hfib{f}b}$ 对所有 $b:B$ 都是可缩的，从而 $\trunc n{\hfib{\varphi}{\pairr{b,v}}}$ 对所有 $(b,v):\sm{b:B} Q(b)$ 也是可缩的。
另一方面，若 $\varphi$ 是 $n$-连通的，则 $\trunc n{\hfib{\varphi}{\pairr{b,v}}}$ 对所有 $(b,v)$ 都是可缩的，因此对任何存在 $v:Q(b)$ 的 $b:B$，$\trunc n{\hfib{f}b}$ 也是可缩的。
最后，由于可缩性是一个命题，只需仅知存在这样的 $v$ 即可。
\end{proof}

如果 $Q$ 不是逐纤维仅知有实例的，\cref{lem:nconnected_postcomp_variation} 的逆命题可能不成立。
例如，如果 $P$ 和 $Q$ 都恒取 $\emptyt$，那么 $\varphi$ 和每个 $g_a$ 都是等价，但 $f$ 可以是任意的。

在另一个方向上，我们有

\begin{lem}\label{prop:nconn_fiber_to_total}
设 $P,Q:A\to\type$ 是类型族，并考虑一个从 $P$ 到 $Q$ 的逐纤维变换\index{fiberwise!transformation}
\begin{equation*}
f:\prd{a:A} \Parens{P(a)\to Q(a)}
\end{equation*}
。
那么诱导映射 $\total f: \sm{a:A}P(a) \to \sm{a:A} Q(a)$ 是 $n$-连通的，当且仅当每个 $f(a)$ 是 $n$-连通的。
\end{lem}

当然，“仅当”方向也是 \cref{lem:nconnected_postcomp_variation} 的一个特例。

\begin{proof}
根据 \cref{fibwise-fiber-total-fiber-equiv}，对每个 $x:A$ 和 $v:Q(x)$，我们有
$\hfib{\total f}{\pairr{x,v}}\eqvsym\hfib{f(x)}v$
因此 $\trunc n{\hfib{\total f}{\pairr{x,v}}}$ 是可缩的，当且仅当
$\trunc n{\hfib{f(x)}v}$ 是可缩的。
\end{proof}

关于连通映射的另一个有用事实是，它们在 $n$-截断上诱导一个等价：

\begin{lem} \label{lem:connected-map-equiv-truncation}
若 $f : A \to B$ 是 $n$-连通的，则它诱导一个等价
$\eqv{\trunc{n}{A}}{\trunc{n}{B}}$。
\end{lem}

\begin{proof}
令 $c$ 为 $f$ 是 $n$-连通的证明。从左到右，我们使用映射 $\trunc{n}{f} : \trunc{n}{A} \to \trunc{n}{B}$。
为了定义从右到左的映射，根据截断的泛性质，只需给出一个映射 $\mathsf{back} : B \to {\trunc{n}{A}}$。
我们可以如下定义这个映射：
\[
\mathsf{back}(y) \defeq \trunc{n}{\proj{1}}{(\proj{1}{(c(y))})}.
\]
根据定义，$c(y)$ 的类型是 $\iscontr(\trunc n {\hfiber{f}y})$，所以它的第一个分量的类型是 $\trunc n{\hfiber{f}y}$，而我们通过投影可以从中得到一个 $\trunc n A$ 的元素。

接下来，我们证明这两个复合映射都是恒等映射。在两个方向上，由于目标都是在 $n$-截断类型中给出一条道路，只需处理构造子 $\tprojf{n}$ 的情形。

在一个方向上，需证对任意 $x:A$，
\[
\trunc{n}{\proj{1}}{(\proj{1}{(c(f(x)))})} = \tproj{n}{x}.
\]
但 $\tproj{n}{(x, \refl{f(x)})} : \trunc n{\hfiber{f}{f(x)}}$，而 $c(f(x))$ 断言这个类型是可缩的，因此
\[
\proj{1}{(c(f(x)))} = \tproj{n}{(x, \refl{})}.
\]
将 $\trunc{n}{\proj{1}}$ 应用于该等式两边即得结论。

在另一个方向上，需证对任意 $y:B$，
\[
\trunc{n}{f}(\trunc{n}{\proj{1}} (\proj{1}{(c(y))})) = \tproj{n}{y}.
\]
$\proj{1}{(c(y))}$ 的类型是 $\trunc n {\hfiber{f}y}$，我们想要的道路本质上是 $\hfiber{f}y$ 的第二个分量，但要确保截断的部分能对应上。

一般地，假设给定 $p:\trunc{n}{\sm{x:A} B(x)}$ 并希望证明
$P(\trunc{n}{\proj{1}{}}(p))$。根据截断归纳法，只需对所有 $a:A$ 和 $b:B(a)$ 证明 $P(\tproj{n}{a})$。
将这一原理应用于当前情形，只需在给定 $a:A$ 和 $b:f (a) = y$ 时证明
\[
\trunc{n}{f}(\tproj{n}{a}) = \tproj{n}{y}
\]
但左边等于 $\tproj{n}{f (a)}$，所以将 $\tprojf{n}$ 应用于 $b$ 两边即得结论。
\end{proof}

有人可能会猜想这个事实完全刻画了 $n$-连通映射，但实际上 $n$-连通性比这稍强一些。
例如，包含映射 $\bfalse:\unit \to\bool$ 在 $(-1)$-截断上诱导一个等价，但它不是满的（即不是 $(-1)$-连通的）。
在 \cref{sec:long-exact-sequence-homotopy-groups} 中我们将看到，一般的差异在于一种类似的额外满射性。

\section{正交分解}
\label{sec:image-factorization}

\index{unique!factorization system|(}%
\index{orthogonal factorization system|(}%
在集合论中，满射和单射构成一个唯一分解系统：每个函数本质上唯一地分解为一个满射后接一个单射。
我们已经看到，满射可以自然地推广到 $n$-连通映射，因此很自然地会问：这些映射是否也参与一个分解系统？
以下是单射的相应推广。

\begin{defn}
  一个函数 $f:A\to B$ 是\define{$n$-截断的（$n$-truncated）}
  \indexdef{n-truncated@$n$-truncated!function}%
  \indexdef{function!n-truncated@$n$-truncated}%
  如果对任意 $b:B$，纤维 $\hfib f b$ 都是 $n$-类型。
\end{defn}

特别地，$f$ 是 $(-2)$-截断的当且仅当它是等价。
当然，$A$ 是 $n$-类型当且仅当 $A\to\unit$ 是 $n$-截断的。
此外，$n$-截断映射也可以像 $n$-类型那样递归地定义。

\begin{lem}\label{thm:modal-mono}
  对任意 $n\ge -2$，函数 $f:A\to B$ 是 $(n+1)$-截断的当且仅当对任意 $x,y:A$，映射 $\apfunc{f}:(x=y) \to (f(x)=f(y))$ 是 $n$-截断的。
  \index{function!embedding}%
  \index{function!injective}%
  特别地，$f$ 是 $(-1)$-截断的当且仅当它是 \cref{sec:mono-surj} 中意义上的嵌入。
\end{lem}

\begin{proof}
  注意对任意 $(x,p),(y,q):\hfib f b$，有
  \begin{align*}
    \big((x,p) = (y,q)\big)
    &= \sm{r:x=y} (p = \apfunc f(r)\ct q)\\
    &= \sm{r:x=y} (\apfunc f (r) = p\ct \opp q)\\
    &= \hfib{\apfunc{f}}{p\ct \opp q}.
  \end{align*}
  因此，$f$ 的任意纤维中的任意道路空间都是 $\apfunc{f}$ 的纤维。
  另一方面，取 $b\defeq f(y)$ 和 $q\defeq \refl{f(y)}$，我们看到 $\apfunc f$ 的任意纤维都是 $f$ 的纤维中的一个道路空间。
  由于 $f$ 是 $\nplusone$-截断的当且仅当其所有纤维的道路空间都是 $n$-类型，结论得证。
\end{proof}

现在我们可以以一种相当明显的方式构造分解。

\begin{defn}\label{defn:modal-image}
设 $f:A\to B$ 是一个函数。$f$ 的\define{$n$-像（$n$-image）}
\indexdef{image}%
\indexdef{image!n-image@$n$-image}%
\indexdef{n-image@$n$-image}%
\indexdef{function!n-image of@$n$-image of}%
定义为
\begin{equation*}
\im_n(f)\defeq \sm{b:B} \trunc n{\hfib{f}b}.
\end{equation*}
当 $n=-1$ 时，我们简记 $\im(f)$ 并称其为 $f$ 的\define{像（image）}。
\end{defn}

\begin{lem}\label{prop:to_image_is_connected}
对任意函数 $f:A\to B$，典范映射 $\tilde{f}:A\to\im_n(f)$ 是 $n$-连通的。
因此，任意函数可分解为一个 $n$-连通映射后接一个 $n$-截断映射。
\end{lem}

\begin{proof}
注意 $A\eqvsym\sm{b:B}\hfib{f}b$。映射 $\tilde{f}$ 是由典范的逐纤维变换
\begin{equation*}
\prd{b:B} \Parens{\hfib{f}b\to\trunc n{\hfib{f}b}}.
\end{equation*}
诱导出的全空间上的映射。
由于每个映射 $\hfib{f}b\to\trunc n{\hfib{f}b}$ 根据 \cref{cor:totrunc-is-connected} 是 $n$-连通的，故由 \cref{prop:nconn_fiber_to_total} 知 $\tilde{f}$ 是 $n$-连通的。
最后，投影 $\proj1:\im_n(f) \to B$ 是 $n$-截断的，因为它的纤维等价于 $f$ 的纤维的 $n$-截断。
\end{proof}

在接下来的引理中，我们将为证明唯一分解定理准备一些工具。

\begin{lem}\label{prop:factor_equiv_fiber}
假设我们有函数的交换图
\begin{equation*}
  \xymatrix{
    {A} \ar[r]^{g_1} \ar[d]_{g_2} &
    {X_1} \ar[d]^{h_1} &
    \\
    {X_2} \ar[r]_{h_2}
    &
    {B}
  }
\end{equation*}
其中$H:h_1\circ g_1\htpy h_2\circ g_2$，且 $g_1$ 与 $g_2$ 是 $n$-连通的，而 $h_1$ 与 $h_2$ 是 $n$-截断的。
则对任意 $b:B$，存在一个等价
\begin{equation*}
E(H,b):\hfib{h_1}b\eqvsym\hfib{h_2}b
\end{equation*}
使得对任意 $a:A$，有一个相等证明
\[\overline{E}(H,a) :  E(H,h_1(g_1(a)))({g_1(a),\refl{h_1(g_1(a))}}) = \pairr{g_2(a),\opp{H(a)}}.\]
\end{lem}

\begin{proof}
任取 $b:B$。则我们有以下等价：
\begin{align}
\hfib{h_1}b
& \eqvsym \sm{w:\hfib{h_1}b} \trunc n{ \hfib{g_1}{\proj1 w}}
\tag{since $g_1$ is $n$-connected}\\
& \eqvsym \Trunc n{\sm{w:\hfib{h_1}b}\hfib{g_1}{\proj1 w}}
\tag{by \cref{thm:refl-over-ntype-base}, since $h_1$ is $n$-truncated}\\
& \eqvsym \trunc n{\hfib{h_1\circ g_1}b}
\tag{by \cref{ex:unstable-octahedron}}
\end{align}
对 $h_2$ 和 $g_2$ 同理。
此外，由于有同伦$H:h_1\circ g_1\htpy h_2\circ g_2$，存在一个明显的等价$\hfib{h_1\circ g_1}b\eqvsym\hfib{h_2\circ g_2}b$。
因此对任意 $b:B$，我们得到
\begin{equation*}
\hfib{h_1}b\eqvsym\hfib{h_2}b
\end{equation*}
通过分析底层的函数，我们得到元素
$\pairr{g_1(a),\refl{h_1(g_1(a))}}$ 在 $E$ 的各组成等价作用下变化的如下表示。
其中有一些相等证明是定义性的，但另一些（在下方用 $=$ 标出）仅是命题性的；将它们组合起来，便得到 $\overline E(H,a)$。
{\allowdisplaybreaks
\begin{align*}
\pairr{g_1(a),\refl{h_1(g_1(a))}} &
    \overset{=}{\mapsto} \Pairr{\pairr{g_1(a),\refl{h_1(g_1(a))}}, \tproj n{ \pairr{a,\refl{g_1(a)}}}}\\
  & \mapsto \tproj n { \pairr{\pairr{g_1(a),\refl{h_1(g_1(a))}}, \pairr{a,\refl{g_1(a)}} }}\\
  & \mapsto \tproj n { \pairr{a,\refl{h_1(g_1(a))}}}\\
  & \overset{=}{\mapsto} \tproj n { \pairr{a,\opp{H(a)}}}\\
  & \mapsto \tproj n { \pairr{\pairr{g_2(a),\opp{H(a)}},\pairr{a,\refl{g_2(a)}}} }\\
  & \mapsto \Pairr{\pairr{g_2(a),\opp{H(a)}}, \tproj n {\pairr{a,\refl{g_2(a)}}} }\\
  & \mapsto \pairr{g_2(a),\opp{H(a)}}
\end{align*}}
第一个等式成立是因为：对于一般的 $b$，映射
\narrowequation{ \hfib{h_1}b \to \sm{w:\hfib{h_1}b} \trunc n{ \hfib{g_1}{\proj1 w}} }
会插入$\trunc n{ \hfib{g_1}{\proj1 w}}$的收缩中心（由"$g_1$ 是 $n$-截断的"这一假设提供）；而在当前讨论的情形中，该类型有明显的元素$\tproj n{ \pairr{a,\refl{g_1(a)}}}$，由可缩性可知它必须等于那个收缩中心。
第二个命题性等式成立是因为等价$\hfib{h_1\circ g_1}b\eqvsym\hfib{h_2\circ g_2}b$将第二个分量与 $\opp{H(a)}$ 拼接，并且我们有$\opp{H(a)} \ct \refl{} = \opp{H(a)}$。
其余等式均为定义性的，读者可以自行验证（假定\cref{ex:unstable-octahedron}有一个合理的解决方案）。
\end{proof}

% The equivalences $E(H,b)$ are such that $E(H^{-1},b)= E(H,b)^{-1}$.

结合\cref{prop:to_image_is_connected,prop:factor_equiv_fiber}，我们得到如下唯一分解结果：

\begin{thm}\label{thm:orth-fact}
对每个 $f:A\to B$，由下式定义的空间 $\fact_n(f)$
\begin{equation*}
\sm{X:\type}{g:A\to X}{h:X\to B} (h\circ g\htpy f)\times\mathsf{conn}_n(g)\times\mathsf{trunc}_n(h)
\end{equation*}
是可缩的。
其收缩中心是来自\cref{prop:to_image_is_connected}的元素
\begin{equation*}
\pairr{\im_n(f),\tilde{f},\proj1,\theta,\varphi,\psi}:\fact_n(f)
\end{equation*}
其中$\theta:\proj1\circ\tilde{f}\htpy f$是典范同伦，$\varphi$ 是
\cref{prop:to_image_is_connected}的证明，而 $\psi$ 是"$\proj1:\im_n(f)\to B$具有 $n$-截断的纤维"这一事实的显然证明。
\end{thm}

\begin{proof}
根据\cref{prop:to_image_is_connected}我们知道 $\fact_n(f)$ 中存在元素，因此只需证明 $\fact_n(f)$ 是一个命题。
假设 $f$ 有两个 $n$-分解
\begin{equation*}
\pairr{X_1,g_1,h_1,H_1,\varphi_1,\psi_1}\qquad\text{and}\qquad\pairr{X_2,g_2,h_2,H_2,\varphi_2,\psi_2}
\end{equation*}
那么我们有点态拼接而成的同伦
\[ H\defeq (\lam{a} H_1(a) \ct H_2^{-1}(a)) \,:\, (h_1\circ g_1\htpy h_2\circ g_2).\]
根据泛等性以及道路与搬运在 $\Sigma$-类型、函数类型和道路类型中的刻画，只需证明以下四点：
\begin{enumerate}
\item 存在一个等价 $e:X_1\eqvsym X_2$，
\item 存在一个同伦$\zeta:e\circ g_1\htpy g_2$，
% \note{Is it easy enough to see that these elements are the various transports?}
\item 存在一个同伦$\eta:h_2\circ e\htpy h_1$，
\item 对任意 $a:A$ 我们有$\opp{\apfunc{h_2}(\zeta(a))} \ct \eta(g_1(a)) \ct H_1(a) = H_2(a)$。
\end{enumerate}
%where $\underline{e}$ is the function underlying the equivalence.
我们按顺序证明这四个断言。
\begin{enumerate}
\item 由\cref{prop:factor_equiv_fiber}，我们有一个逐纤维等价
% \note{It could be a nice exercise for the book to show
% that if $f_1:A_1\to B$ and $f_2:A_2\to B$ have equivalent fibers, then $A_1\eqvsym A_2$}.
\begin{equation*}
E(H) : \prd{b:B} \eqv{\hfib{h_1}b}{\hfib{h_2}b}.
\end{equation*}
这诱导了全空间的一个等价，亦即我们有
\begin{equation*}
\eqvspaced{\Parens{\sm{b:B} \hfib{h_1}b}}{\Parens{\sm{b:B}\hfib{h_2}b}}.
\end{equation*}
当然，由\cref{thm:total-space-of-the-fibers}我们还有等价$X_1\eqvsym\sm{b:B}\hfib{h_1}b$与$X_2\eqvsym\sm{b:B}
\hfib{h_2}b$。
由此我们得到等价 $e:X_1\eqvsym X_2$；读者可以验证 $e$ 的底层函数由下式给出：
\begin{equation*}
e(x) \jdeq \proj1(E(H,h_1(x))(x,\refl{h_1(x)})).
\end{equation*}
\item 由\cref{prop:factor_equiv_fiber}，我们可以选取
  $\zeta(a) \defeq \apfunc{\proj1}(\overline E(H,a)) : e(g_1(a)) = g_2(a)$。
  \label{item:orth-fact-2}
\item 对每个 $x:X_1$，我们有
\begin{equation*}
\proj2(E(H,h_1(x))({x,\refl{h_1(x)}})) :h_2(e(x))= h_1(x),
\end{equation*}
这给出了同伦$\eta:h_2\circ e\htpy h_1$。
\item 根据纤维中道路的刻画（\cref{lem:hfib}），来自\cref{prop:factor_equiv_fiber}的道路 $\overline E(H,a)$ 给出
  $\eta(g_1(a)) = \apfunc{h_2}(\zeta(a)) \ct \opp{H(a)}$。
  代入 $H$ 的定义并整理道路后，即得所求的等式。\qedhere
\end{enumerate}
\end{proof}

% I can't make sense of this, and it doesn't seem necessary
%
% \begin{cor}
% A function $f:A\to B$ is $n$-connected if and only if
% \begin{equation*}
% \prd C\prd{g:\modalfunc(B\to C)} \iscontr\big(\sm{h:\modalfunc(B\to C)}\underline{h}\circ
% f\htpy\underline{g}\circ f\big).
% \end{equation*}
% \end{cor}

通过标准论证，我们得到如下的正交性原理。

\begin{thm}
  设 $e:A\to B$ 是 $n$-连通的，且 $m:C\to D$ 是 $n$-截断的。
  则映射
  \[ \varphi: (B\to C) \;\to\; \sm{h:A\to C}{k:B\to D} (m\circ h \htpy k \circ e) \]
  是一个等价。
\end{thm}

\begin{proof}[证明概要]
  对于陪域中的任意 $(h,k,H)$，令 $h = h_2 \circ h_1$ 且 $k = k_2 \circ k_1$，其中 $h_1$ 与 $k_1$ 是 $n$-连通的，$h_2$ 与 $k_2$ 是 $n$-截断的。
  那么$f = (m\circ h_2) \circ h_1$与$f = k_2 \circ (k_1\circ e)$都是 $m \circ h = k\circ e$ 的 $n$-分解。
  因此，二者之间存在唯一的等价。
  由此不难得出$\hfib\varphi{(h,k,H)}$是可缩的（尽管略显繁琐）。
\end{proof}

\index{orthogonal factorization system|)}%
\index{unique!factorization system|)}%

最后，我们证明像在拉回下是稳定的。
\index{image!stability under pullback}
\index{factorization!stability under pullback}

\begin{lem}\label{lem:hfiber_wrt_pullback}
假设方块
\begin{equation*}
  \vcenter{\xymatrix{
      A\ar[r]\ar[d]_f &
      C\ar[d]^g\\
      B\ar[r]_-h &
      D
      }}
\end{equation*}
是一个拉回方块，并设 $b:B$。那么$\hfib{f}b\eqvsym\hfib{g}{h(b)}$。
\end{lem}

\begin{proof}
这可由拉回的粘贴性质（\cref{ex:pullback-pasting}）得出，因为下图中的类型 $X$
\begin{equation*}
  \vcenter{\xymatrix{
      X\ar[r]\ar[d] &
      A\ar[r]\ar[d]_f &
      C\ar[d]^g\\
      \unit\ar[r]_b &
      B\ar[r]_h &
      D
      }}
\end{equation*}
是左边方块的拉回当且仅当它是整个外层矩形的拉回，而 $\hfib{f}b$ 是左边方块的拉回，$\hfib{g}{h(b)}$ 是整个外层矩形的拉回。
\end{proof}

\begin{thm}\label{thm:stable-images}
\index{stability!of images under pullback}%
考虑函数 $f:A\to B$, $g:C\to D$ 以及图表
\begin{equation*}
  \vcenter{\xymatrix{
      A\ar[r]\ar[d]_{\tilde{f}_n} &
      C\ar[d]^{\tilde{g}_n}\\
      \im_n(f)\ar[r]\ar[d]_{\proj1} &
      \im_n(g)\ar[d]^{\proj1}\\
      B\ar[r]_h &
      D
      }}
\end{equation*}
如果外层矩形是拉回，那么底部的方块也是拉回（从而由\cref{ex:pullback-pasting}，顶部的方块也是拉回）。因此，像在拉回下是稳定的。
\end{thm}

\begin{proof}
假设外层方块是拉回，我们有等价
\begin{align*}
B\times_D\im_n(g) & \jdeq \sm{b:B}{w:\im_n(g)} h(b)=\proj1 w\\
& \eqvsym \sm{b:B}{d:D}{w:\trunc n{\hfib{g}d}} h(b)= d\\
& \eqvsym \sm{b:B} \trunc n{\hfib{g}{h(b)}}\\
& \eqvsym \sm{b:B} \trunc n{\hfib{f}b} &&
\text{(by \cref{lem:hfiber_wrt_pullback})}\\
& \equiv \im_n(f). && \qedhere
\end{align*}
\end{proof}

\index{n-type@$n$-type|)}%

\section{模态}
\label{sec:modalities}

\index{modality|(}

几乎所有关于 $n$-类型和连通性的理论都可以在更一般的框架下进行。
本节的内容在本书后续部分不会用到。

对于推广 $n$-类型理论，我们最初的想法可能是将\cref{thm:trunc-reflective}作为一个定义。

\begin{defn}\label{defn:reflective-subuniverse}
  一个\define{反射子宇宙（reflective subuniverse）}
  \indexdef{reflective!subuniverse}%
  \indexdef{subuniverse, reflective}%
  是一个谓词 $P:\type\to\prop$，满足对于每个类型 $A:\type$，我们都有一个类型 $\reflect A$ 使得 $P(\reflect A)$ 成立，以及一个映射
  $\project_A:A\to\reflect A$，并具有如下性质：对于每个满足 $P(B)$ 的类型 $B:\type$，下面的映射是一个等价：
  \[\function{(\reflect A\to{}B)}{(A\to{}B)}{f}{f\circ\project_A}.\]
\end{defn}

我们记$\P \defeq \setof{A:\type | P(A)}$，因此 $A:\P$ 表示 $A:\type$ 且 $P(A)$ 成立。
我们也将上述映射的拟逆记为 $\rec{\modal}$。
记号 $\reflect$ 可能看起来有点奇怪，但稍后我们会明白它的意义。

对于任意反射子宇宙，我们可以用通常的方式证明范畴论中关于反射子范畴的所有熟知结论。例如，我们有：

\begin{itemize}
\item 一个类型 $A$ 属于 $\P$ 当且仅当 $\project_A:A\to\reflect A$ 是一个等价。
\item $\P$ 在收缩核下封闭。
  特别地，只要 $\project_A$ 有一个收缩，$A$ 就属于 $\P$。
\item 运算 $\reflect$ 在某种合适的、同伦相干的意义下是一个函子；我们可以在需要的任何高阶层次上给出精确表述。
\item $\P$ 中的类型在所有极限（如积和拉回）下封闭。
  特别地，对于任意 $A:\P$ 和 $x,y:A$，恒等类型 $(x=_A y)$ 也属于 $\P$，因为它是两个函数 $\unit\to A$ 的拉回。
\item $\P$ 中的余极限可以通过将 $\reflect$ 作用于类型的普通余极限来构造。
\end{itemize}

重要的是，对积的封闭性也推广到"无穷积"，即依值函数类型。

\begin{thm}\label{thm:reflsubunv-forall}
  如果 $B:A\to\P$ 是反射子宇宙 \P 中任意一族类型，那么 $\prd{x:A} B(x)$ 也属于 \P。
\end{thm}

\begin{proof}
  对于任意 $x:A$，考虑由$\mathsf{ev}_x(f) \defeq f(x)$定义的函数$\mathsf{ev}_x : (\prd{x:A} B(x)) \to B(x)$。
  因为 $B(x)$ 属于 $P$，这个函数可以延拓为一个函数
  \[ \rec{\modal}(\mathsf{ev}_x) : \reflect\Parens{\prd{x:A} B(x)} \to B(x). \]
  因此我们可以通过$h(z)(x) \defeq \rec{\modal}(\mathsf{ev}_x)(z)$定义$h:\reflect(\prd{x:A} B(x)) \to \prd{x:A} B(x)$。
  那么 $h$ 就是$\project_{\prd{x:A} B(x)}$的一个收缩，从而${\prd{x:A} B(x)}$属于 $\P$。
\end{proof}

特别地，如果 $B:\P$ 且 $A$ 是任意类型，那么 $(A\to B)$ 属于 \P。
用范畴论的语言来说，这意味着任何反射子宇宙都是一个\define{指数理想（exponential ideal）}。
\indexdef{exponential ideal}%
进而，由标准论证可知反射子保持有限积。

\begin{cor}\label{cor:trunc_prod}
  对于任意类型 $A$ 和 $B$ 以及任意反射子宇宙，诱导的映射$\reflect(A\times B) \to \reflect(A) \times \reflect(B)$是一个等价。
\end{cor}

\begin{proof}
  只需证明 $\reflect(A) \times \reflect(B)$ 具有与 $\reflect(A\times B)$ 相同的泛性质。
  根据上述关于 $\P$ 中类型在极限下封闭的说明，$\reflect(A) \times \reflect(B)$ 属于 $\P$。
  现在令 $C:\P$；我们有
  \begin{align*}
    (\reflect(A) \times \reflect(B) \to C)
    &= (\reflect(A) \to (\reflect(B) \to C))\\
    &= (\reflect(A) \to (B \to C))\\
    &= (A \to (B \to C))\\
    &= (A \times B \to C)
  \end{align*}
  这里利用了 $\reflect(B)$ 和 $\reflect(A)$ 的泛性质，以及 $B\to C$ 属于 $\P$ 的事实（因为 $C$ 属于 $\P$）。
  可以直接验证这个等价是通过与 $\mreturn_A \times \mreturn_B$ 复合得到的，符合要求。
\end{proof}

类型的每个反射子范畴都自动成为指数理想，并且具有保持积的反射子，这可能显得奇怪。然而，在经典集合范畴中，出于同样的原因，情况也是如此。只是这一事实通常不被提及，因为经典的集合范畴——与同伦类型范畴相反——并没有许多有趣的反射子范畴。

$n$-类型的两个基本性质不被一般的反射子宇宙所具有：\cref{thm:ntypes-sigma}（在 $\Sigma$-类型下封闭）和\cref{thm:truncn-ind}（截断归纳）。然而，这两个性质的类似命题彼此等价。

\begin{thm}\label{thm:modal-char}
  对于一个反射子宇宙 \P，以下两条逻辑等价。
  \begin{enumerate}
  \item 如果 $A:\P$ 且 $B:A\to \P$，那么 $\sm{x:A} B(x)$ 属于 \P。\label{item:mchr1}
  \item 对于每个 $A:\type$，类型族 $B:\reflect A\to\P$，以及映射$g:\prd{a:A} B(\project(a))$，存在$f:\prd{z:\reflect A} B(z)$使得对所有 $a:A$ 有$f(\project(a)) = g(a)$。\label{item:mchr2}
  \end{enumerate}
\end{thm}

\begin{proof}
  假设\ref{item:mchr1}。
  那么在\ref{item:mchr2}所述的情形中，类型$\sm{z:\reflect A} B(z)$属于 $\P$，并且我们有由$g'(a)\defeq (\project(a),g(a))$定义的$g':A\to \sm{z:\reflect A} B(z)$。
  于是，我们得到$\rec{\modal}(g'):\reflect A \to \sm{z:\reflect A} B(z)$，使得$\rec{\modal}(g')(\project(a)) = (\project(a),g(a))$成立。

  现在考虑函数$\proj1 \circ \rec{\modal}(g') : \reflect A \to \reflect A$ 与 $\idfunc[\reflect A]$。
  根据假设，它们与 $\project$ 预复合后相等。
  因此，由 $\reflect$ 的泛性质，它们已经相等，即对所有 $z$ 我们有$p_z:\proj1(\rec{\modal}(g')(z)) = z$。
  现在我们可以定义
  %
  \narrowequation{f(z) \defeq \trans{p_z}{\proj2(\rec{\modal}(g')(z))},}
  %
  利用定义\ref{defn:reflective-subuniverse}给出的等价的伴随性质，可以证明
  %
  \narrowequation{\rec{\modal}(g')(\project(a)) = (\project(a),g(a))}
  %
  的第一分量等于 $p_{\project(a)}$。因此，它的第二分量给出所需的$f(\project(a)) = g(a)$。

  反过来，假设\ref{item:mchr2}，且 $A:\P$ 以及 $B:A\to\P$。
  令 $h$ 为复合
  \[ \reflect\Parens{\sm{x:A} B(x)} \xrightarrow{\reflect(\proj1)} \reflect A \xrightarrow{\opp{(\project_A)}} A. \]
  那么对于$z:\sm{x:A} B(x)$，我们有
  \begin{align*}
    h(\project(z)) &= \opp\project(\reflect(\proj1)(\project(z)))\\
    &= \opp\project(\project(\proj1(z)))\\
    &= \proj1(z).
  \end{align*}
  记这条道路为 $p_z$。
  现在如果我们将$C:\reflect(\sm{x:A} B(x)) \to \type$定义为$C(w) \defeq B(h(w))$，则有
  \[ g \defeq \lam{z} \trans{p_z}{\proj2(z)} \;:\; \prd{z:\sm{x:A} B(x)} C(\project(z)). \]
  于是，由假设可得
  %
  \narrowequation{f:\prd{w:\reflect(\sm{x:A}B(x))} C(w)}
  %
  使得$f(\project(z)) = g(z)$成立。
  $h$ 和 $f$ 共同给出一个函数
  %
  \narrowequation{k:\reflect(\sm{x:A}B(x)) \to \sm{x:A}B(x)}
  %
  由$k(w) \defeq (h(w),f(w))$定义，而 $p_z$ 与相等$f(\project(z)) = g(z)$表明 $k$ 是$\project_{\sm{x:A}B(x)}$的一个收缩。
  因此，$\sm{x:A}B(x)$ 属于 \P。
\end{proof}

注意这与\cref{sec:htpy-inductive}中的讨论相似。
\index{recursion principle!for a modality}%
\index{induction principle!for a modality}%
\index{uniqueness!principle, propositional!for a modality}%
反射子宇宙的反射子的泛性质类似于带有唯一性性质的递归原理，而\cref{thm:modal-char}\ref{item:mchr2}则类似于相应的归纳原理。
与\cref{sec:htpy-inductive}中不同，由于我们只能消去到属于 $\P$ 的类型中，二者在这里并不等价。
\cref{thm:modal-char} 的条件~\ref{item:mchr1} 正是用来弥合这一脱节的。

当然，毫不意外的是，如果我们有归纳原理，就可以推出递归原理。
只要允许消去到道路类型，我们也能推出其唯一性性质。
这引出了以下定义。
注意，任何反射子宇宙都可以由操作 $\reflect:\type\to\type$ 和函数 $\project_A:A\to \reflect A$ 来刻画，因为$P(A) = \isequiv(\project_A)$。

\begin{defn}\label{defn:modality}
一个 \define{模态（modality）}
\indexdef{modality}
是一个操作 $\modal:\type\to\type$，并附带以下数据：
\begin{enumerate}
\item 对每个类型 $A$，有一个函数 $\mreturn^\modal_A:A\to\modal(A)$。\label{item:modal1}
\item 对每个 $A:\type$ 以及每个类型族 $B:\modal(A)\to\type$，有一个函数\label{item:modal2}
\begin{equation*}
\ind{\modal}:\Parens{\prd{a:A}\modal(B(\mreturn^\modal_A(a)))}\to\prd{z:\modal(A)}\modal(B(z)).
\end{equation*}
\item 对每个 $f:\prd{a:A}\modal(B(\mreturn^\modal_A(a)))$，有一条道路 $\ind\modal(f)(\mreturn^\modal_A(a)) = f(a)$。\label{item:modal3}
\item 对任意 $z,z':\modal(A)$，函数 $\mreturn^\modal_{z=z'} : (z=z') \to \modal(z=z')$ 是一个等价。\label{item:modal4}
\end{enumerate}
如果 $\mreturn^\modal_A:A\to\modal(A)$ 是一个等价，则称 $A$ 对于 $\modal$ 是 \define{模态的}
\indexdef{modal!type}%
\indexdef{type!modal}%
，并记
\begin{equation}
  \modaltype\defeq\setof{X:\type | X \text{ is $\modal$-modal} }\label{eq:modaltype}
\end{equation}
为模态类型的类型。
\end{defn}

条件~\ref{item:modal2} 和~\ref{item:modal3} 与 \cref{thm:modal-char}\ref{item:mchr2} 非常相似，但改用 $\modal B(z)$ 来表述，而非假设 $B$ 取值于 $\P$。
这使得我们可以纯粹用操作 $\modal$ 来陈述条件，而无需预先给定谓词 $P:\type\to\prop$。
（这并非完全令人满意，因为在条款~\ref{item:modal4}中我们仍然不太隐讳地用到了 $P$。
我们不知道~\ref{item:modal4} 是否能从~\ref{item:modal1}--\ref{item:modal3} 推出。）
然而，\cref{thm:modal-char}\ref{item:mchr2} 这个看起来更强的性质可以从 \cref{defn:modality}\ref{item:modal2} 及~\ref{item:modal3} 推出，因为对于任何 $C:\modal A \to \modaltype$ 都有 $C(z) \eqvsym \modal C(z)$，并且我们可以借助这个等价来回转换。

\index{universal!property!of a modality}%
如同其他归纳原理一样，这蕴含了一个泛性质。

\begin{thm}\label{prop:lv_n_deptype_sec_equiv_by_precomp}
设 $A$ 是一个类型，且 $B:\modal(A)\to\modaltype$。则函数
\begin{equation*}
(\blank\circ \mreturn^\modal_A) : \Parens{\prd{z:\modal(A)}B(z)} \to \Parens{\prd{a:A}B(\mreturn^\modal_A(a))}
\end{equation*}
是一个等价。
\end{thm}

\begin{proof}
由定义，操作 $\ind{\modal}$ 是 $(\blank\circ \mreturn^\modal_A)$ 的一个右逆。
因此，我们只需对每个 $s:\prd{z:\modal(A)}B(z)$ 找到一个同伦
\begin{equation*}
\prd{z:\modal(A)}s(z)= \ind{\modal}(s\circ \mreturn^\modal_A)(z)
\end{equation*}
来证明它同时也是一个左逆。
根据假设，每个 $B(z)$ 都是模态的，因而每个类型 $s(z)= R^\modal_X(s\circ \mreturn^\modal_A)(z)$
也是模态的。
于是，只需找到一个类型为
\begin{equation*}
\prd{a:A}s(\mreturn^\modal_A(a))= \ind{\modal}(s\circ \mreturn^\modal_A)(\mreturn^\modal_A(a))
\end{equation*}
的函数即可，而这可由 \cref{defn:modality}\ref{item:modal3} 推出。
\end{proof}

特别地，对每个类型 $A$ 和每个模态类型 $B$，我们都有一个等价 $(\modal A\to B)\eqvsym (A\to B)$。

\begin{cor}
  对于任意模态 $\modal$，$\modal$-模态类型构成一个满足 \cref{thm:modal-char} 的等价条件的反射子宇宙。
\end{cor}

因此，模态可以等同于对 $\Sigma$ 类型封闭的反射子宇宙。
名称 \emph{模态（modality）} 当然来源于 \emph{模态逻辑（modal logic）}\index{modal!logic}，后者研究的是一种逻辑，在其中我们可以形成诸如"可能 $A$"（通常写作 $\diamond A$）或"必然 $A$"（通常写作 $\Box A$）这类陈述。
符号 $\modal$ 通常用于表示任意的模态算子\index{modal!operator}。% (rather than a specific one such as $\diamond$ or $\Box$).
在命题即类型的原则下，模态逻辑意义上的模态对应着对\emph{类型}的一种操作，而 \cref{defn:modality} 似乎是定义此类操作的一种合理候选方案。
（更准确地说，我们或许应该称这些模态为\emph{幂等的、单子的（idempotent, monadic）}模态；见注记。）
\index{idempotent!modality}%
如 \cref{subsec:when-trunc} 所述，我们通常可以使用副词\index{adverb}来非正式地谈论这些模态，例如用"仅仅"\index{merely}表示命题截断，用"纯粹"\index{purely}表示恒等模态
\index{identity!modality}%
\index{modality!identity}%
（即由 $\modal A \defeq A$ 定义的那个模态）。

对于任意模态 $\modal$，我们称映射 $f:A\to B$ 是 \define{$\modal$-连通的（$\modal$-connected）}
\indexdef{function!.circle-connected@$\modal$-connected}%
\indexdef{.circle-connected function@$\modal$-connected function}%
，若对所有 $b:B$，$\modal(\hfib f b)$ 都是可缩的；称其为 \define{$\modal$-截断的（$\modal$-truncated）}
\indexdef{function!.circle-truncated@$\modal$-truncated}%
\indexdef{.circle-truncated function@$\modal$-truncated function}%
，若对所有 $b:B$，$\hfib f b$ 都是模态的。
所有来自 \cref{sec:connectivity,sec:image-factorization} 的理论，只要不涉及关联不同 $n$ 值的 $n$-类型，都可以逐字逐句地适用于这一一般情形。
\index{orthogonal factorization system}%
\index{unique!factorization system}%
特别地，我们有一个正交分解系统。

一类不包含 $n$-截断的重要模态是\emph{左正合（left exact）}模态：即那些其函子 $\modal$ 既保持拉回又保持有限积的模态。
\index{topology!Lawvere-Tierney}%
它们是初等意象\index{topos}论中"Lawvere–Tierney\index{Lawvere}\index{Tierney} 拓扑"的范畴化，并在高阶范畴语义中对应于子 $(\infty,1)$-意象。
\index{.infinity1-topos@$(\infty,1)$-topos}%
然而，这已超出了本书的范围。

除 $n$-截断外，其他模态的一些具体例子可以在习题中找到。

\index{modality|)}

\sectionNotes

同伦 $n$-类型这一概念在经典同伦论中由来已久。
是 Voevodsky 认识到，这个概念可以在同伦类型论中递归地定义，从可缩性出发即可。

\index{axiom!Streicher's Axiom K}%
"Axiom K"这一性质由 Thomas Streicher 如此命名，是因为它是继 J 之后相等类型的一个性质——后者是相等类型消去子的传统名称。
\cref{thm:hedberg} 归功于 Hedberg~\cite{hedberg1998coherence}；\cite{krausgeneralizations} 则包含了更多信息与推广。

$n$-连通空间与函数的概念在同伦论中同样经典，不过如前所述，我们对函数连通性的编号与经典编号相差 1。
由此产生的分解系统的重要性，已被 Rezk、Lurie 及其他学者在高阶意象理论中的近期工作所强调。%
\index{.infinity1-topos@$(\infty,1)$-topos}
特别地，本章的结果应与~\cite[\S6.5.1]{lurie:higher-topoi} 相比较。
在 \cref{sec:freudenthal} 中，$n$-连通映射的理论将对我们证明 Freudenthal 纬悬定理起到关键作用。

对\emph{简单}类型论中的模态算子\index{modal!operator}已有大量研究；例如可见~\cite{modalTT}。而在依值类型论的框架下，\cite{ab:bracket-types} 将命题截断（$(-1)$-截断）这一特例作为模态算子\index{modal!operator}来处理。这里所展示的发展极大地扩展并推广了上述工作，同时也借鉴了意象理论\index{topos}中的思想。

一般来说，模态算子\index{modal!operator}至少可以分为两种：一种是像 $\diamond$（“可能”）那样满足 $A\Rightarrow \diamond A$ 的，另一种是像 $\Box$（“必然”）那样满足 $\Box A \Rightarrow A$ 的。
当它们还是\emph{幂等的}（即 $\diamond A = \diamond{\diamond A}$ 或 $\Box A = \Box{\Box A}$）时，前者可以等同于反射子范畴（或等价地，幂等单子），而后者则等同于余反射子范畴（或幂等余单子）。
\index{monad}
\index{comonad}
然而，在依值类型论中，处理余单子那一类要棘手一些，因为它们在拉回之下更少能保持稳定，从而无法被解释为宇宙 \UU 上的运算。
虽然有时有办法绕过这一点（例如参见\cite{QGFTinCHoTT12}），但为了简单起见，这里我们只讨论单子那一类。

在计算方面，单子（从而模态\index{modality}）被用来为函数式编程中的计算效应建模~\cite{Moggi89}。%
\index{programming}%
\index{computational effect}
如果一个计算的执行不会导致副作用（例如在屏幕上打印消息、播放音乐或通过互联网发送数据），就称该计算是\emph{纯}的。
存在像 Haskell\index{Haskell} 这样的“纯函数式”编程语言，在技术上，其中只能编写纯函数：副作用是通过对输出类型应用“单子”来表示的。
例如，类型为 $\mathsf{Int}\to\mathsf{Int}$ 的函数是纯的，而类型为 $\mathsf{Int}\to \mathsf{IO}(\mathsf{Int})$ 的函数在计算结果的过程中可能会执行输入和输出；运算 $\mathsf{IO}$ 就是一个单子。
\index{purely}%
（这就是我们将副词 “purely” 用于恒等单子的原因，因为后者在计算上对应于没有副作用的纯函数。）
本章考虑的所有模态都是幂等的，而函数式编程中使用的模态很少是幂等的，但两者的思想仍然密切相关。

\sectionExercises

\begin{ex}\label{ex:all-types-sets}\
  \begin{enumerate}
    \item 利用 \cref{thm:h-set-refrel-in-paths-sets} 证明：如果对每个类型 $A$ 都有 $\brck{A}\to A$，那么每个类型都是集合。
    \item 证明：如果每个满射函数都（纯粹地）分裂，即如果对每个 $f:A\to B$ 都有
    %
    \narrowequation{\prd{b:B}\brck{\hfib{f}{b}}\to\prd{b:B}\hfib{f}{b}},
    %
    那么每个类型都是集合。
  \end{enumerate}
\end{ex}

\begin{ex}\label{ex:s2-colim-unit}
  对于本题，我们考虑以下关于余极限的一般概念。
  定义一个\define{图（graph）}\indexdef{graph} $\Gamma$，它由一个类型 $\Gamma_0$ 和一个族 $\Gamma_1 : \Gamma_0 \to \Gamma_0 \to \UU$ 组成。
  在一个图 $\Gamma$ 上的\define{图表（diagram）}\index{diagram}（指类型的图表）包含一个族 $F:\Gamma_0 \to \UU$，以及对每对 $x,y:\Gamma_0$，一个函数 $F_{x,y}:\Gamma_1(x,y) \to F(x) \to F(y)$。
  该图表的\define{余极限（colimit）}\index{colimit!of types} 是由以下构造子生成的高阶归纳类型 $\colim(F)$：
  \begin{itemize}
  \item 对每个 $x:\Gamma_0$，一个函数 $\inc_x:F(x) \to \colim(F)$，以及
  \item 对每个 $x,y:\Gamma_0$、每个 $\gamma:\Gamma_1(x,y)$ 以及每个 $a:F(x)$，一条道路 $\inc_y(F_{x,y}(\gamma,a)) = \inc_x(a)$。
  \end{itemize}
  也存在更一般的余极限（例如参见\cref{ex:s2-colim-unit-2}），但这对于许多目的而言已经足够。
  \begin{enumerate}
  \item 给出一个图 $\Gamma$，使得 $\Gamma$-图表的余极限可以等同于 \cref{sec:colimits} 中定义的推出。
    换言之，每个伸展应该诱导出一个 $\Gamma$-图表，其余极限就是该伸展的推出。
  \item 给出一个图 $\Gamma$ 及其上的图表 $F$，使得对所有 $x$ 有 $F(x)=\unit$，但同时有 $\colim(F)=\Sn^2$。
    注意 $\unit$ 是一个 $(-2)$-类型，而 $\Sn^2$ 预期不是任何有限 $n$ 对应的 $n$-类型。
    另见 \cref{ex:s2-colim-unit-2}。
  \end{enumerate}
\end{ex}

\begin{ex}\label{ex:ntypes-closed-under-wtypes}
  证明：如果 $A$ 是一个 $n$-类型且 $B:A\to \ntype{n}$ 是一个 $n$-类型的族，其中 $n\ge -1$，那么 $W$-类型 $\wtype{a:A} B(a)$（参见 \cref{sec:w-types}）也是一个 $n$-类型。
\end{ex}

\begin{ex}\label{ex:connected-pointed-all-section-retraction}
  利用 \cref{prop:nconn_fiber_to_total} 将 \cref{thm:connected-pointed} 推广到任意的截面-收缩对。
\end{ex}

\begin{ex}\label{ex:ntype-from-nconn-const}
  证明 \cref{thm:nconn-to-ntype-const} 的特征刻画在另一个方向上也成立：$B$ 是 $n$-类型，当且仅当每个从 $n$-连通类型到 $B$ 的映射都是常值的。
  理想情况下，你的证明应该能像在 \cref{sec:modalities} 中那样，适用于任何模态。
\end{ex}

\begin{ex}\label{ex:connectivity-inductively}
  证明：对于 $n\ge -1$，一个类型 $A$ 是 $n$-连通的，当且仅当它仅知有实例，并且对于所有 $a,b:A$，类型 $\id[A]ab$ 是 $(n-1)$-连通的。
  因此，由于每个类型都是 $(-2)$-连通的，类型的 $n$-连通性可以仅使用命题截断来归纳地定义。
  （特别地，$A$ 是 $0$-连通的当且仅当 $\brck{A}$ 且 $\prd{a,b:A} \brck{a=b}$。）
\end{ex}

\begin{ex}\label{ex:lemnm}
  \indexdef{excluded middle!LEMnm@$\LEM{n,m}$}%
  对于 $-1\le n,m \le\infty$，记 $\LEM{n,m}$ 为如下命题
  \[ \prd{A:\ntype{n}} \trunc m{A + \neg A},\]
  这里将 $\ntype{\infty} \defeq \type$，且 $\trunc{\infty}{X}\defeq X$。
  证明：
  \begin{enumerate}
  \item 若 $n=-1$ 或 $m=-1$，则 $\LEM{n,m}$ 等价于 \cref{sec:intuitionism} 中的 $\LEM{}$。
  \item 若 $n\ge 0$ 且 $m\ge 0$，则 $\LEM{n,m}$ 与泛等性矛盾。
  \end{enumerate}
\end{ex}

\begin{ex}\label{ex:acnm}
  \indexdef{axiom!of choice!ACnm@$\choice{n,m}$}%
  对于 $-1\le n,m\le\infty$，记 $\choice{n,m}$ 为如下命题
  \[ \prd{X:\set}{Y:X\to\ntype{n}}
  \Parens{\prd{x:X} \trunc m{Y(x)}}
  \to
  \Trunc m{\prd{x:X} Y(x)},
  \]
  约定同 \cref{ex:lemnm}。
  于是 $\choice{0,-1}$ 就是 \cref{sec:axiom-choice} 中的选择公理，
  而 $\choice{\infty,\infty}$ 就是恒等函数。
  （如果我们按 \eqref{eq:ac} 而非 \eqref{eq:epis-split} 的方式类似地定义 $\choice{n,m}$，
  那么 $\choice{\infty,\infty}$ 就会类似于 \cref{thm:ttac}。）
  已知 $\choice{\infty,-1}$ 与泛等性一致，因为它在 Voevodsky 的单纯集模型中成立。
  \begin{enumerate}
  \item 不使用泛等性，证明对任意 $m$，$\LEM{n,\infty}$ 蕴含 $\choice{n,m}$。
    （另一方面，在 \cref{subsec:emacinsets} 中我们将证明 $\choice{}=\choice{0,-1}$ 蕴含 $\LEM{}=\LEM{-1,-1}$。）
  \item 显然，如果 $k\le n$，则 $\choice{n,m}\Rightarrow \choice{k,m}$。
    那么这些原理 $\choice{n,m}$ 之间还有其他蕴含关系吗？
    对于任意 $m\ge 0$ 和任意 $n$，$\choice{n,m}$ 与泛等性一致吗？
    （这些都是开放问题。）\index{open!problem}
  \end{enumerate}
\end{ex}

\begin{ex}\label{ex:acnm-surjset}
  证明 $\choice{n,-1}$ 蕴含以下结论：对任意 $n$-类型 $A$，都仅存在一个集合 $B$ 和一个满射 $B\to A$。
\end{ex}

\begin{ex}\label{ex:acconn}
  \define{$n$-连通选择公理（$n$-connected axiom of choice）}
  \indexdef{n-connected@$n$-connected!axiom of choice}%
  \indexdef{axiom!of choice!n-connected@$n$-connected}%
  指的是如下命题：
  \begin{quote}
    若 $X$ 是一个集合，且 $Y:X\to \type$ 是一个类型族，使得每个 $Y(x)$ 都是 $n$-连通的，则 $\prd{x:X} Y(x)$ 也是 $n$-连通的。
  \end{quote}
  注意 $(-1)$-连通选择公理就是 \cref{ex:acnm} 中的 $\choice{\infty,-1}$。
  \begin{enumerate}
  \item 证明对所有 $n\ge -1$，$(-1)$-连通选择公理蕴含 $n$-连通选择公理。
  \item $n$-连通选择公理与原理 $\choice{n,m}$ 之间还有其他的蕴含关系吗？
    （这是一个开放问题。）\index{open!problem}
  \end{enumerate}
\end{ex}

\begin{ex}\label{ex:n-truncation-not-left-exact}
  证明对任意 $n\ge -1$，$n$-截断模态都不是左正合的。
  也就是说，请给出一个拉回，使得该模态未能保持它。
\end{ex}

\begin{ex}\label{ex:double-negation-modality}
  证明 $X\mapsto (\neg\neg X)$ 是一个模态。\index{modal!operator}%
\end{ex}

\begin{ex}\label{ex:prop-modalities}
  设 $P$ 是一个命题。
  \begin{enumerate}
  \item 证明 $X\mapsto (P\to X)$ 是一个左正合模态。
    这称为与 $P$ 关联的\define{开模态（open modality）}
    \indexdef{open!modality}%
    \indexdef{modality!open}%
    。
  \item 证明 $X\mapsto P*X$ 是一个左正合模态，其中 $*$ 表示联结（见 \cref{sec:colimits}）。
    这称为与 $P$ 关联的\define{闭模态（closed modality）}
    \indexdef{closed!modality}%
    \indexdef{modality!closed}%
    。
  \end{enumerate}
\end{ex}

\begin{ex}\label{ex:f-local-type}
  设 $f:A\to B$ 是一个映射；如果 $(\blank\circ f):(B\to Z) \to (A\to Z)$ 是一个等价，则称类型 $Z$ 是\define{$f$-局部的（$f$-local）}
  \indexdef{f-local type@$f$-local type}%
  \indexdef{type!f-local@$f$-local}%
  。
  \begin{enumerate}
  \item 证明 $f$-局部类型构成一个反射子宇宙。
    你需要使用一个高阶归纳类型来定义反射子（局部化）。
  \item 证明如果 $B=\unit$，那么这个子宇宙是一个模态。
  \end{enumerate}
\end{ex}

\begin{ex}\label{ex:trunc-spokes-no-hub}
  证明：与 \cref{rmk:spokes-no-hub} 不同，我们可以等价地将 $\trunc nA$ 定义为由一个函数 $\tprojf n : A \to \trunc n A$ 以及对于每个 $r:\Sn^{n+1} \to \trunc n A$ 和每个 $x:\Sn^{n+1}$ 的一条路径 $s_r(x):r(x) = r(\base)$ 所生成。
\end{ex}

\begin{ex}\label{ex:s2-colim-unit-2}
  本练习中，我们考虑一种比 \cref{ex:s2-colim-unit} 中稍复杂一些的余极限概念。
  定义\define{带复合的图（graph with composition）}\indexdef{graph!with composition} $\Gamma$ 为如 \cref{ex:s2-colim-unit} 所述的图，并附带对每个 $x,y,z:\Gamma_0$ 的一个函数 $\Gamma_1(y,z) \to \Gamma_1(x,y) \to \Gamma_1(x,z)$，记作 $\delta\mapsto\gamma \mapsto \delta \circ \gamma$。
  （例如，任何如 \cref{cha:category-theory} 所述的预范畴都是带复合的图。）
  带复合的图 $\Gamma$ 上的\define{图表（diagram）}\index{diagram} $F$ 由底层图上的图表，加上对每个 $x,y,z:\Gamma_0$、$\gamma:\Gamma_1(x,y)$ 和 $\delta:\Gamma_1(y,z)$ 的一个同伦 $\cmp_{x,y,z}(\delta,\gamma) : F_{y,z}(\delta) \circ F_{x,y}(\gamma) \htpy F_{x,z}(\delta\circ\gamma)$ 构成。
  此类图表的\define{余极限（colimit）}\index{colimit!of types} 是由以下构造子生成的高阶归纳类型 $\colim(F)$：
  \begin{itemize}
  \item 对每个 $x:\Gamma_0$，有一个函数 $\inc_x:F(x) \to \colim(F)$，
  \item 对每个 $x,y:\Gamma_0$、$\gamma:\Gamma_1(x,y)$ 和 $a:F(x)$，有一条道路 $\glue_{x,y}(\gamma,a) : \inc_y(F_{x,y}(\gamma,a)) = \inc_x(a)$，以及
  \item 对每个 $x,y,z:\Gamma_0$、$\gamma:\Gamma_1(x,y)$、$\delta:\Gamma_1(y,z)$ 和 $a:F(x)$，有一条道路
    \[ \ap{\inc_z}{\cmp_{x,y,z}(\delta,\gamma,a)} \ct \glue_{x,z}(\delta\circ \gamma,a) = \glue_{y,z}(\delta,F_{x,y}(\gamma,a)) \ct \glue_{x,y}(\gamma,a). \]
  \end{itemize}
  （这是对完全同伦论意义上的图表和余极限概念的一种“二阶近似”；后者在所有更高层级上均应包含此类“相干性路径”。
  在类型论中定义此类概念是一个重要的公开问题。）

  试构造一个带复合的图 $\Gamma$，使得 $\Gamma_0$ 是集合且每个类型 $\Gamma_1(x,y)$ 都是命题，并构造其上的图表 $F$，使得对所有 $x$ 均有 $F(x)=\unit$，且 $\colim(F)=\Sn^2$。
\end{ex}

\begin{ex}\label{ex:fiber-map-not-conn}
  对比 \cref{lem:nconnected_postcomp_variation,prop:nconn_fiber_to_total}，人们可能会猜想：若 $f:A\to B$ 是 $n$-连通的，且 $g:\prd{a:A} P(a) \to Q(f(a))$ 诱导出 $n$-连通映射 $\Parens{\sm{a:A} P(a)} \to \Parens{\sm{b:B} Q(b)}$，则 $g$ 是逐纤维 $n$-连通的。
  请给出一个反例以说明此猜想不成立。
  （事实上，推广到模态时，这一性质刻画了左正合模态；参见 \cref{ex:prop-modalities}。）
\end{ex}

\begin{ex}\label{ex:is-conn-trunc-functor}
  证明：若 $f : A \to B$ 是 $n$-连通的，则 $\trunc kf : \trunc kA \to \trunc kB$ 也是 $n$-连通的。
\end{ex}

\begin{ex}\label{ex:categorical-connectedness}
  我们称一个类型 $A$ 是\define{范畴意义下连通的（categorically connected）}
  \indexdef{connected!categorically}，如果对于任意类型 $B, C$，由
  \begin{align*}
    e_{A,B,C}(\inl(g)) &\defeq \lam{x} \inl(g(x)),\\
    e_{A,B,C}(\inr(g)) &\defeq \lam{x} \inr(g(x))
  \end{align*}
  定义的典范映射 $e_{A, B, C}:((A\to B) + (A\to C)) \to (A \to B + C)$ 是一个等价。
  \begin{enumerate}
  \item 证明任何连通的类型都是范畴意义下连通的。
  \item 证明所有范畴意义下连通的类型都是连通的，当且仅当 $\LEM{}$ 成立。（提示：考虑 $A \defeq \Sigma P$，其中 $\neg \neg P$ 成立。）
  \end{enumerate}
\end{ex}

%%% Local Variables:
%%% mode: latex
%%% TeX-master: "hott-online"
%%% End: